% concatenated from main.tex
\documentclass[12pt]{amsart}

%\usepackage{amssymb,amsfonts,commath,subcaption,caption,amsmath, mathtools}
%\usepackage{graphicx, bbm, tipa,tikz-cd, faktor}
\usepackage{mathtools}
\usepackage{titlesec,caption, subcaption}
\usepackage{tikz-cd}
\usepackage{amsmath,amsthm,amsfonts,amscd,amssymb,amsopn,enumerate,enumitem,amssymb}
%\usepackage{eucal}
\usepackage{mathrsfs,xcolor,hyperref,stmaryrd}
\usepackage{wasysym}
%\usepackage{pxfonts}
\usepackage[all]{xy}
\usepackage{fullpage}
\usepackage{verbatim}
\usepackage{colonequals}
\usepackage[new]{old-arrows}
\usepackage{tocloft}
%\hypersetup{hidelinks}
\titleformat{name=\section}{}{\thetitle.}{0.8em}{\centering\scshape}
\titleformat{name=\subsection}[runin]{}{\thetitle.}{0.8em}{\bfseries}[.]
\titleformat{name=\subsubsection}[runin]{}{\thetitle.}{0.8em}{\bfseries}[.]


\DeclareRobustCommand\longtwoheadrightarrow
     {\relbar\joinrel\twoheadrightarrow}

\theoremstyle{plain}
\newtheorem{thm}{Theorem}[subsection]
\newtheorem*{thm*}{Theorem}
\newtheorem{cor}[thm]{Corollary}
\newtheorem*{assumption}{(*)}
\newtheorem{lem}[thm]{Lemma}
\newtheorem{conj}[thm]{Conjecture}
\newtheorem*{conj*}{Conjecture}
\newtheorem*{rem*}{Remark}
\newtheorem{fact}[thm]{Fact}
\newtheorem*{fact*}{Fact}
\newtheorem*{prop*}{Proposition}

\theoremstyle{definition}
\newtheorem{defn}[thm]{Definition}
\newtheorem*{defn*}{Definition}
\newtheorem{defns}[thm]{Definitions}
\newtheorem*{ack}{Acknowledgement}
\newtheorem{con}[thm]{Construction}
\newtheorem{rem}[thm]{Remark}
\newtheorem{exmps}[thm]{Examples}
\newtheorem{notn}[thm]{Notation}
\newtheorem{notns}[thm]{Notations}
\newtheorem{addm}[thm]{Addendum}
\newtheorem{exer}[thm]{Exercise}
\newtheorem{prob}[thm]{Problem}
\newtheorem{prop}[thm]{Proposition}
\newtheorem{theorem}[thm]{Theorem}
\newtheorem{lemma}[thm]{Lemma}
\theoremstyle{remark}

\newtheorem{rems}[thm]{Remarks}
\newtheorem{warn}[thm]{Warning}
\newtheorem{sch}[thm]{Scholium}
\newtheorem{exmp}[thm]{Example}


\renewcommand{\baselinestretch}{1.3}

\newcommand{\medcap}{\mathbin{\scalebox{1.5}{\ensuremath{\cap}}}}
\newcommand{\medcup}{\mathbin{\scalebox{1.5}{\ensuremath{\cup}}}}
\newcommand{\lege}[2]{\left(\frac{#1}{#2}\right)}
\newcommand{\loc}[1]{#1_{\mathfrak{p}}}
\newcommand{\locat}[2]{#1_{\mathfrak{#2}}}
\renewcommand{\thefootnote}{}
%\setlength{\parindent}{0pt}
%\parskip0.5em

%\addtolength{\oddsidemargin}{-.875in}
%\addtolength{\evensidemargin}{-.875in}
%\addtolength{\textwidth}{1.60in}

%\addtolength{\topmargin}{-1in}
%\addtolength{\textheight}{1.60in}
%of 
\begin{document}
	
	
\title{On the pro-modularity in the residually reducible case for some totally real fields}
\author{Xinyao Zhang}
\maketitle

\begin{abstract}
	In this article, we study the relation between the universal deformation rings and big Hecke algebras in the residually reducible case. Following the strategy of Skinner-Wiles and Pan's proof of the Fontaine-Mazur conjecture, we prove a pro-modularity result. Based on this result, we also give a conditional big $R=\mathbb{T}$ theorem over some totally real fields, which is a generalization of Deo's result.
\end{abstract}


\footnotetext{2020 Mathematics Subject Classification: Primary - 11F80; Secondary - 11F85}
\footnotetext{Keywords: pseudo-representation; modularity}

\section{Introduction}

Let $\bar{\rho}_0: \operatorname{Gal}(\bar{\mathbb{Q}}/\mathbb{Q}) \to \operatorname{GL}_2(\mathbb{F})$ be a continuous odd representation, where $\mathbb{F}$ is finite field with characteristic $p \ge 3$. Let $R_{\bar{\rho}_0}$ be the universal deformation ring of $\bar{\rho}_0$. It is conjectured by Gouv\^ea (\cite{gouvea2006arithmetic}) that $R_{\bar{\rho}_0}$ is isomorphic to some $p$-adic big Hecke algebra. When $\bar{\rho}_0$ is absolutely irreducible, this question is studied by B\"ockle (\cite{bockle2001density}) and Allen (\cite{allen2019automorphic}). Recently, Deo has explored the reducible case in \cite{deo2023density}. In this article, we try to generalize his results to some totally real fields.

More precisely, suppose that $\bar{\rho}_0^{\textnormal{ss}} =1 \oplus \bar{\chi}$ for some character $\bar{\chi}$ unramified outside $Np$, where $N$ is a positive integer. Let $R_{\bar{\rho}_x}$ be the universal deformation ring of $\rho_x$, where $0 \ne x \in H^1(G_{\mathbb{Q}, Np}, \mathbb{F}(\bar{\chi}))$ and $\rho_x$ is the unique reducible Galois representation determined by $x$ such that $\rho_x^{\textnormal{ss}}=1 \oplus \bar{\chi}$. Let $R_{\bar{\rho}_0}^{\textnormal{pd}}$ be the universal pseudo-deformation ring of $\bar{\rho}_0^{\textnormal{ss}}$. Deo proves the following results.

%1) \cite[Corollary 3.5]{deo2023density} The ring $R_{\bar{\rho}_x}$ is a local complete intersection ring of Krull dimension $4$.

\begin{thm}{\cite[Theorem A, Theorem 4.5]{deo2023density}}
	Let $N_0$ be the tame Artin conductor of $\bar{\rho}_0$. Suppose $N_0|N$, $\bar{\chi}|_{G_{\mathbb{Q}_p}} \ne \mathbf{1}, \omega_p^{-1}$, $\operatorname{dim}_{\mathbb{F}} H^1(G_{\mathbb{Q}, Np}, \mathbb{F}(\bar{\chi}))=1$ and $p \nmid \phi(N)$. Then
	
	1) There is an isomorphism $ (R_{\bar{\rho}_0}^{\textnormal{pd}})^{\operatorname{red}} \cong T(N)_{\bar{\rho}_0}$, where $T(N)_{\bar{\rho}_0}$ is the localization of the big Hecke algebra $T(N)$ (of level $N$) at the maximal ideal $\mathfrak{m}_0$ determined by $\bar{\rho}_0$.
	
	2) There exists an isomorphism $ R_{\bar{\rho}_x} \twoheadrightarrow T(N)_{\bar{\rho}_0}$.
\end{thm} 
 
In his proof, a crucial assumption is the cyclicity of $H^1(G_{\mathbb{Q}, Np}, \mathbb{F}(\bar{\chi}))$. Under this assumption, Deo uses the fact that the natural map $R_{\bar{\rho}_0}^{\textnormal{pd}} \to R_{\bar{\rho}_x} $ is actually a surjection proved by Kisin \cite[Corollary 1.4.4]{kisin2009fontaine} in the proof of 1), and constructs a surjection $R_{\bar{\rho}_x} \twoheadrightarrow T(N)_{\bar{\rho}_0} $ induced by $R_{\bar{\rho}_0}^{\textnormal{pd}} \twoheadrightarrow T(N)_{\bar{\rho}_0}$ by the theory of Generalized Matrix Algebra in the proof of 2). However, such a cyclicity assumption is absurd when we consider a totally real field $F$ instead of $\mathbb{Q}$ due to the global characteristic formula.

To generalize Deo's results to some totally real fields, we need a new strategy, which is inspired by the proof of the Fontaine-Mazur conjecture in the residually reducible case from Skinner-Wiles (\cite{skinner1999residually}) and Pan (\cite{pan2022fontaine}).
Specifically, we study the pro-modularity of the universal deformation rings.

Now we state our main results.

 Let $p>3$ be a rational prime. Now we assume that $F$ is an abelian totally real field of even degree such that $p$ splits completely in $F$. Let $\Sigma$ be a finite set of finite places of $F$ containing $\Sigma_p $ consisting of all the places $v \mid p$ . Let $\mathcal{O}$ be the ring of integers of some local field $E$ with a uniformizer $\pi$ and residue field $\mathbb{F}$, and $\chi: \operatorname{Gal}(\bar{F}/F) \to \mathcal{O}^\times$ a totally odd character unramified outside $\Sigma$, and $\bar{\chi}=\chi ~\operatorname{mod} ~\pi$. Let $F_{\Sigma}$ be the maximal extension of $F$ unramified outside $\Sigma$ and all infinite places. Let $\epsilon, \omega_p$ be the $p$-adic, $\operatorname{mod}~p$ cyclotomic character, respectively.
 
%and $\mathbb{F}$ be a finite field of characteristic $p$

For an element $0 \ne x \in H^1(F_{\Sigma}/F, \mathbb{F}(\bar{\chi}^{-1}))$, let $\rho_x$ be the unique non-split reducible Galois representation determined by $x$ such that $\rho_x^{\textnormal{ss}}=1 \oplus \bar{\chi}$. Let $R_x^\textnormal{aux}$ be the universal deformation ring of $\rho_x$ of auxiliary type (see Definition \ref{auxiliary}).
%\begin{thm} [Proposition \ref{dim aux}]
	%Keep notations and assumptions as above. Assume that $F$ is an abelian totally real field. Further assume that $$[F:\mathbb{Q}]=\operatorname{dim}_{\mathbb{F}}(H^1(F_{\Sigma}/F, \mathbb{F}(\bar{\chi}^{-1}))) = \operatorname{dim}_{\mathbb{F}} (H^1(F_{\Sigma}/F, \mathbb{F}(\bar{\chi}))).$$
	%Then $R_x^\textnormal{aux}$ is a local complete intersection ring of Krull dimension $1+2[F:\mathbb{Q}]$.
%\end{thm}
  Let $R_\textnormal{aux}^\textnormal{ps} $ be the universal pseudo-deformation ring of the residual pseudo-representation $1+\bar{\chi}$ with fixed determinant $\chi$. Let $\mathbb{T}_{\xi}$ be the big $p$-adic Hecke algebra (see Definition \ref{big}), and assume that $ \bar{\rho}_0^{\textnormal{ss}}$ determines a maximal ideal $\mathfrak{m}_{\xi}$. For a prime $\mathfrak{p}$ of $R_\textnormal{aux}^\textnormal{ps} $ (resp. $R_x^\textnormal{aux}$), we say that it is pro-modular if it comes from a prime of $(\mathbb{T}_{\xi})_{\mathfrak{m}_{\xi}}$ (see Definition \ref{nice} and Definition \ref{pro-modular} for the explicit definitions). 

\begin{thm} [Theorem \ref{pro}, Theorem \ref{MINIMAL} \& Theorem \ref{special case}] \label{main} Keep notations as above. For any place $v \in \Sigma \setminus \Sigma_p$, assume that $p \mid \operatorname{Nm}(v)-1$. Assume that the character $\chi$ satisfies the following conditions:
	
	1) $\chi$ is unramified outside $ \Sigma_p$,
	
	2) $\chi(\operatorname{Frob}_v) \equiv 1~\operatorname{mod}~\pi$ for  
	$v \in \Sigma \setminus \Sigma_p$.
	
	3) $\bar{\chi}$ can be extended to a character of $G_{\mathbb{Q}}$.
	
	4) $\bar{\chi} \mid_{G_{F_v}} \ne 1, \omega_p^{-1}$ for any $v \mid p$.
	
	5) $\chi \mid_{G_{F_v}}$ is de Rham for any $v \mid p$.
	
	(1) If $[F:\mathbb{Q}] \ge 7|\Sigma \setminus \Sigma_p|+4$, then for any irreducible component of $R_\textnormal{aux}^\textnormal{ps} $ of dimension at least $1+2[F:\mathbb{Q}]$, it is pro-modular of dimension $1+2[F:\mathbb{Q}]$.
	
   (2) If $[F:\mathbb{Q}] \ge \operatorname{max}\{1+\frac{1}{2}\operatorname{dim}_{\mathbb{F}}(H^1(F_{\Sigma}/F, \mathbb{F}(\bar{\chi}^{-1}))), 7|\Sigma \setminus \Sigma_p|+4\}$. Then for any non-zero cocycle $x \in H^1(F_{\Sigma}/F, \mathbb{F}(\bar{\chi}^{-1}))$, any prime ideal of $ R_x^\textnormal{aux}$ is pro-modular. Further, $R_x^\textnormal{aux}$ is a local complete intersection ring of Krull dimension $1+2[F:\mathbb{Q}]$.
\end{thm}

	To prove such a theorem, we follow Pan's strategy in \cite{pan2022fontaine}, and a key tool is the patching argument (Theorem \ref{R=T}) at a nice prime (see Definition \ref{nice}). Using Theorem \ref{R=T}, we only need to find a nice prime in each irreducible component of large dimension in the universal pseudo-deformation ring. It can be found successfully if the degree $[F: \mathbb{Q}]$ is large enough using Corollary \ref{innovation'}.

\begin{rem}
	One may ask how we can use these pro-modularity results.
	
	The first part of Theorem \ref{main}, i.e., the pro-modularity of the universal pseudo-deformation ring, is actually a complement of the first step in Pan's proof of the Fontaine-Mazur conjecture in the residually reducible case (see \cite[page 1032]{pan2022fontaine}) so that after an abelian base change, the conjecture is a direct consequence of (1) of Theorem \ref{main} and \cite[Corollary 3.5.10]{pan2022fontaine}, although the full pro-modularity is not needed in his proof. In fact, Pan's strategy is to find a ``potentially nice prime'', i.e., it becomes nice  after some soluble base change,  and it is much easier to construct. Then he can identify some irreducible components of the universal (pseudo-)deformation ring as in the big Hecke algebra, which is enough for his purpose.
	
	 However, in the missing case of the Fontaine-Mazur conjecture in \cite{pan2022fontaine}, i.e., $p=3$ and $\bar{\chi}|_{G_{F_v}}=\omega_p^{-1}$ for any $v \mid p$, to study the pro-modularity of some irreducible component of the universal (pseudo-)deformation ring, we need to know the pro-modularity of some other ones, which means we actually need a (potentially) big $R=\mathbb{T}$ theorem in this case. More precisely, the ``dihedral condition'' (the third one in (2) of Definition \ref{nice}) sometimes obstructs the existence of the potentially nice primes, whereas the strategy of finding nice primes in this article, especially the application of Corollary \ref{innovation'}, still works so that we just need to find pro-modular primes rather than nice primes (see Proposition \ref{promodular to nice}). In other words, our method in this paper, combining some classical patching arguments, can be able to deal with some subtle cases of the conjecture, especially the case $p=3$ and $ \bar{\chi}=\omega_3$. Therefore, it allows us to drop Pan's second assumption of \cite[Theorem 1.0.4]{pan2022fontaine} and hence concludes the two-dimensional Fontaine-Mazur conjecture in the regular case for odd primes. See Remark \ref{adv} and \cite[Section 5]{ZHANG2024} for more details.
	 
	 We expect that our method can be generalized to prove more modularity or automorphy lifting theorems in the residually reducible case.
\end{rem}

\begin{rem}
	One may also ask whether we can find some examples such that our theorems may apply.
	
	Our assumptions on the global  character $\chi$ roughly state that it ``comes from" the determinant of a continuous two-dimensional geometric Galois representation $\rho$ of $G_{\mathbb{Q}}$. In other words, starting from such a character $\chi=\operatorname{det} \rho$, using Grunwald-Wang's theorem (see \cite[Theorem 5, page 80]{artin}), we can find a totally real abelian number field $F$ satisfying all of our assumptions of (1) of Theorem \ref{main}. In other words, our pro-modularity result can be viewed as a potential big $R=\mathbb{T}$ theorem for general abelian totally real fields and global characters. This is also the way of applying this theorem to prove the Fontaine-Mazur conjecture.
\end{rem}

	\begin{rem}
		We point out that compared to Deo's results, our pro-modularity result on the universal pseudo-deformation ring does not need the restrictive assumption about the dimension of the cohomology group.
		
		Our pro-modularity result on the universal deformation ring is an analogue to Deo's result \cite[Theorem 4.5]{deo2023density}. It is related to the dimension of the cohomology group since we need to bound the reducible locus of the deformation ring as in \cite{skinner1999residually}.
	\end{rem}
	


%As applications to our pro-modularity arguments, we have the following results.

%\begin{thm} [Theorem \ref{dimT}, Theorem \ref{dimofR}]
	%Keep notations as above.
	
	%(1) If $[F:\mathbb{Q}] \ge 7|\Sigma \setminus \Sigma_p|+2 $, then $(\mathbb{T}_{\xi})_{\mathfrak{m}_{\xi}}$ is equidimensional of dimension $1+2[F: \mathbb{Q}]$.
	
	%(2) With the same assumption of (1), for every irreducible component of $R_\textnormal{aux}^\textnormal{ps} $ of dimension at least $1+2[F:\mathbb{Q}]$, it is exactly of dimension $1+2[F:\mathbb{Q}]$.
	
	%(3) For any non-zero element $x \in H^1(F_{\Sigma}/F, \mathbb{F}(\bar{\chi}^{-1}))$, if $$[F:\mathbb{Q}] \ge \operatorname{max}\{1+\frac{1}{2}\operatorname{dim}_{\mathbb{F}}(H^1(F_{\Sigma}/F, \mathbb{F}(\bar{\chi}^{-1}))), 7|\Sigma \setminus \Sigma_p|+2\},$$
	%then $R_x^\textnormal{aux}$ is a local complete intersection ring of Krull dimension $1+2[F:\mathbb{Q}]$.
	
%\end{thm}

\subsection{Notations and conventions}
Throughout the paper, $F$ is a totally real field and $\mathbb{A}_F$ is the adele ring. For a finite place $v$ of $F$, we write $F_v$ for the completion of $F$ at $v$, $\mathcal{O}_{F_v}$ for its ring of integers, $\pi_v$ for a uniformizer and $k(v)$ for the residue field. We denote the absolute Galois group and the inertia group of $F_v$ by $G_{F_v}$ and $I_v$, respectively. For a rational prime $l$, we denote $\mathcal{O}_F \otimes_\mathbb{Z} \mathbb{Z}_l$ by $\mathcal{O}_{F,l}$. We write $\operatorname{Nm}(v)$ for the norm of $v$, and write $\operatorname{Frob}_v$ for the geometric Frobenius element. 

We fix a finite field $\mathbb{F}$ of characteristic $p>3$. Let $\mathcal{O}$ be a ring of integers of some local field $E$ (a finite extension of $\mathbb{Q}_p$) with a uniformizer $\pi$ and residue field $\mathbb{F}$. Let $\mathscr{C}_{\mathcal{O}}^{f}$ be the category of Artinian local $\mathcal{O}$-algebras with residue field $\mathbb{F}$, and $\mathscr{C}_{\mathcal{O}}$ be the category of topological local $\mathcal{O}$-algebras which are isomorphic to inverse limits of objects of $\mathscr{C}_{\mathcal{O}}^{f}$. For a universal (pseudo-)deformation ring, we mean a universal object in $\mathscr{C}_{\mathcal{O}}$ which pro-represents the functor from $\mathscr{C}_{\mathcal{O}}^{f}$ to the category of sets sending $R$ to the set of (pseudo-)deformations over $R$.

For a CNL ring $R$, we mean that $R$ is a complete Noetherian local ring with maximal ideal $\mathfrak{m}_R$. If $\mathfrak{p}$ is a prime ideal of $R$, we will denote its residue field by $k(\mathfrak{p})$. We write $R_{\mathfrak{p}}$ for the localization at $\mathfrak{p}$, and $\widehat{R_{\mathfrak{p}}}$ for the $\mathfrak{p}$-adic completion. We write $R^{\operatorname{red}}$ for the maximal reduced quotient of $R$. By $M_i(R)$, we mean the ring of $i \times i$ matrices over $R$.

We write $\epsilon, \omega_p$ for the $p$-adic, $\operatorname{mod}~p$ cyclotomic character, respectively.

\begin{ack}
	The author would like to express appreciation to his supervisor Professor Takeshi Saito for his careful reading,  pointing out some mathematical errors and many helpful suggestions and comments.
\end{ack}

\section{Pseudo-representations}

Pseudo-representations are the main tools to study the modularity of Galois representations in the residually reducible case. We recommend the article \cite{bellaiche2009families} to readers for a general theory of pseudo-representations.

\subsection{Two-dimensional pseudo-representations}
In this section, we will give some results on 2-dimensional pseudo-representations. Our main references are \cite{skinner1999residually} and \cite{pan2022fontaine}. 


\begin{defn} \label{pseudo}
	For a profinite group $G$ and a topological commutative ring $R$ in which 2 is invertible, a 2-dimensional \textit{pseudo-representation} is a continuous function $T: G \to R$ such that 
	
	1) $T(1)=2$,
	
	2) there exists an order 2 element $c \in G$ such that $T(c)=0$,
	
	3) $T(\sigma \tau)=T(\tau \sigma)$ for all $ \sigma, \tau \in G$,
	
	4) $T(\gamma \delta \eta)+T(\gamma \eta \delta)-T(\gamma \eta)T(\delta)-T(\eta \delta)T(\gamma)-T(\delta \gamma)T(\eta)+T(\gamma)T(\delta)T(\eta)=0,$
	for any $\delta, \gamma,\eta\in G$.
	
	The \textit{determinant} $\operatorname{det}(T)$ of a pseudo-representation $T$ is a character (using 4)) defined by $$ \operatorname{det}(T): G \to R^\times, ~
	\operatorname{det}(T)(\sigma)=\frac{1}{2}(T(\sigma)^2-T(\sigma^2)), \sigma \in G.$$
	
	
	
	For a pseudo-representation $T$ and $\sigma, \tau \in G$, define:
	
	1) $a(\sigma)=\frac{1}{2}(T(c\sigma)+T(\sigma))$,
	
	2) $d(\sigma)=\frac{1}{2}(-T(c\sigma)+T(\sigma))$,
	
	3) $y(\sigma, \tau)=a(\sigma \tau)-a(\sigma)a(\tau).$
\end{defn}


If $\rho: G \to \operatorname{GL}_2(R)$ is a continuous representation with an element $c \in G$ satisfying $\rho(c)=\begin{pmatrix}
	1 & ~\\
	~& -1
\end{pmatrix}
$, then $\operatorname{tr}(\rho)$ is a pseudo-representation. More explicitly, if $\rho(\sigma)=\begin{pmatrix}
	a_\sigma & b_\sigma\\
	c_\sigma& d_\sigma
\end{pmatrix}$, then $a(\sigma)=a_\sigma$, $d(\sigma)=d_\sigma$ are as defined above and $y(\sigma, \tau)=b_\sigma c_\tau$.

Pan \cite[2.1.4]{pan2022fontaine} verifies that the continuous functions $\{a,d,y\}$ satisfy the following equations.

(1) $y(\sigma, \tau)=d(\tau\sigma)-d(\tau)d(\sigma)$.

(2) $y(\sigma\tau, \delta)=a(\sigma)y(\tau, \delta)+y(\sigma, \delta)d(\tau)$.

(3) $y(\sigma, \tau\delta)=a(\delta)y(\sigma, \tau)+y(\sigma, \delta)d(\tau)$.

(4) $y(\alpha, \beta)y(\sigma, \tau)=y(\alpha, \tau)y(\sigma, \beta)$.\\

From the last equation, we can conclude the following two easy facts.

\begin{prop}\label{u}
	Let $R$ be an integral domain, and $T: G \to R$ be a pseudo-representation for a profinite group $G$. Let $\alpha, \beta $ be two elements in $ G$.
	
	1) If $y(\alpha, \beta) =0$ in $R$, then either $y(\alpha, \tau)=0$ or $y(\tau, \beta)=0$ for any $\tau \in G$.
	
	2) Assume that the set $J_{\beta} =\{\tau \in G : y(\beta, \tau)\ne 0\}$ (resp. $J'_{\beta} =\{\tau \in G : y(\tau, \beta) \ne 0\}$) is non-empty. For $\tau \in J_{\beta}$ (resp. $J'_{\beta}$), let $u(\alpha, \beta)(\tau)= \frac{y(\alpha, \tau)}{y(\beta, \tau)}$ (resp. $u'(\alpha, \beta)(\tau)=\frac{y(\tau, \alpha)}{y(\tau, \beta)}$) be an element in the fraction field of $R$. Then we have $u(\alpha, \beta)(\tau)=u(\alpha, \beta)(\tau')$ (resp. $u'(\alpha, \beta)(\tau)=u'(\alpha, \beta)(\tau')$)  for any $\tau, \tau' \in J_\beta$ (resp. $J'_{\beta}$). (Afterwards, we will use notations $u(\alpha, \beta)$ and $u'(\alpha, \beta)$ for simplicity.)
\end{prop}

\begin{defn}\label{construction}\cite[2.1.5]{pan2022fontaine}
	Assume that $R$  is either a field or a discrete valuation ring. For a pseudo-representation $T: G\to R$, a representation $\rho$ associated to $T$ is in the following form.
	
	1) $ \rho(\sigma)=
	\begin{pmatrix}
		a(\sigma) & ~\\
		~& d(\sigma)
	\end{pmatrix}$, if all $y(\sigma, \tau)=0$. We call this case reducible.
	
	2) Choose $\sigma_0, \tau_0$ such that $\frac{y(\sigma, \tau)}{y(\sigma_0, \tau_0)} \in R$ for any $ \sigma, \tau$, if $y(\sigma, \tau) \ne 0 $ for some $ \sigma, \tau$. Define $ \rho(\sigma)=
	\begin{pmatrix}
		a(\sigma) & \frac{y(\sigma, \tau_0)}{y(\sigma_0, \tau_0)}\\
		y(\sigma_0, \sigma)& d(\sigma)
	\end{pmatrix}$.  We call this case irreducible.
\end{defn}

\begin{rem}
	We can check that $\rho$ is actually absolutely irreducible, and is unique up to conjugation if $R$ is a field in the irreducible case.
\end{rem}

The following results are essential to our pro-modularity arguments.

\begin{prop}\label{innovation}
	Let $R$ be a CNL domain with maximal ideal $\mathfrak{m}_R$ and residue field $\mathbb{F}$. Let $G$ be a profinite group and $T: G \to R$ be a pseudo-representation of $G$. Assume that $T~\operatorname{mod}~\mathfrak{m}_R$ is reducible. Let $\alpha$ be an element of $G $ and write $I_{\alpha}$ for the ideal of $R$ generated by $\{y(\alpha, \beta), \beta \in G\}$. Assume that any minimal prime belonging to $I_{\alpha}$ (following \cite[page 52]{MR242802}) is of height at least $2$.
	
	1) For any $\sigma \in G$, $u(\sigma, \alpha)$ (defined in Proposition \ref{u}) is well-defined and integral over $R$.
	
	2) If $\alpha'$ is an another element of $G $ such that any minimal prime belonging to $I_{\alpha'}$ of $R$ generated by $\{y(\alpha', \beta), \beta \in G\}$ is of height at least $2$, then $u(\alpha, \alpha')$ is a unit in the normal closure of $R$.
\end{prop}


\begin{proof} 
    Let $B$ be the normal closure of $R$, and $i: R \hookrightarrow B$ be the natural inclusion. As $R$ is a CNL domain, by \cite[Lemma 15.106.3]{stacks-project}, $R$ is unibranch. Thus, $B$ is also a local domain.
    
     We first recall a useful result in commutative algebra.
	
	\begin{lemma}\label{commutative algebra}
		 Let $\mathfrak{p}$ be a prime ideal of $B$ of height $d$. Then $i^{-1}(\mathfrak{p})$ is also of height $d$ in $R$.
	\end{lemma}
	
	\begin{proof}
		This is \cite[Proposition 9.2]{eisenbud2013commutative}.
		
		%We first assume that the $i^{-1}(\mathfrak{p})$ is of height at least $d+1$. As $R$ is a CNL domain, by \cite[Proposition 8.2.29]{liu2006algebraic}, we know that $R$ is a Nagata ring. Then the natural injection $R \hookrightarrow B$ is finite, and thus $\operatorname{dim} R=\operatorname{dim}B$. By \cite[Theorem 8.2.39]{liu2006algebraic}, we know that $R$ is universally catenary, and so is $B$. Hence, we have $\operatorname{dim} B/\mathfrak{p} = \operatorname{dim} B-d$ as $B$ is local. By the incomparability (\cite[Corollary 4.18]{eisenbud2013commutative}), we know that $i^{-1}(\mathfrak{p})$ is of dimension at least $\operatorname{dim} B-d$. From our assumption, we have $\operatorname{dim} R \ge d+1+\operatorname{dim} B-d = \operatorname{dim} B+1$, which is impossible.
		
		%Now assume that $i^{-1}(\mathfrak{p})$ is of height $d' < d$. As $R$ is local and catenary, it is of dimension $\operatorname{dim} R-d'$. By going up property (\cite[Lemma 10.36.22]{stacks-project}), we know that $\mathfrak{p}$ is of dimension at least $\operatorname{dim} R-d'$. As $\operatorname{dim} R=\operatorname{dim}B$, we get a contradiction.
	\end{proof}
	
	Now we return to the proof of our proposition. As $T~\operatorname{mod}~\mathfrak{m}_R$ is reducible, we know that for all $\alpha, \beta \in G$, we have $y(\alpha, \beta) \in \mathfrak{m}_R$.
	
	1) By Lemma \ref{commutative algebra}, we know that there does not exist a prime of $B$ containing $\{y(\alpha, \beta), \beta \in \operatorname{Gal}(F_\Sigma/F)\}$ of height at most $1$. 
	As $B$ is normal , we have $\cap_{\mathfrak{p}_{B}} B_{\mathfrak{p}_{B}}=B $, where $ \mathfrak{p}_{B}$ ranges over all height one primes of $B$. To prove our result, for any height one prime $ \mathfrak{p}_{B}$, we just need to find an elemnt $\beta \in G$ satisfying $\frac{y(\sigma, \beta)}{y(\alpha, \beta)} \in B_{\mathfrak{p}_{B}}$, and it is immediate from our discussion. This proves 1).
	
	2) By 1), we know that both $u(\alpha', \alpha)$ and $u(\alpha, \alpha')$ are well-defined and integral over $R$. As $u(\alpha', \alpha)u(\alpha, \alpha')=1 $, we know that $u(\alpha, \alpha')$ is a unit in $B$. 
\end{proof}

\begin{cor}\label{innovation'}
	Keep notations as in the statement of Proposition \ref{innovation}. Let $S$ be a finite subset (not necessarily a group) of $G$ . Then there exist a partition of $S = S_1 \amalg S_2$, a positive integer $n>1$ satisfying $(n, p)=1$ and a CNL domain $R'$ satisfying the following conditions.
	
	a) $R'$ is a quotient of $R$.
	
	b) For any $\theta \in S_1$, we have $y(\theta, \alpha)=0$ for any $\alpha \in G$ in $R'$. For any $\theta', \theta'' \in S_2$, we have $y(\theta', \alpha)^n=y(\theta'', \alpha)^n$ for any $\alpha \in G$ in $R'$.
	
	c) For any $\alpha \in G,~\theta \in S_2$, either $ y(\theta, \beta)=0$ for any $\beta \in G$ in $R'$ or $u(\alpha, \theta)$ is well-defined and integral over $R'$.
	
	d) We have $\operatorname{dim} R' \ge \operatorname{dim} R - |S|$. If $S_2$ is not empty, then further $\operatorname{dim} R' \ge \operatorname{dim} R - |S|+1$.
	
\end{cor}

\begin{proof}
	We prove this result by induction on $|S|$.
	
	If $|S|=\varnothing$, then our result is just trivial.
	
	Now we assume that our result holds for a non-negative integer $k$ and $|S|=k+1$.
	
	First, we check whether there exists an element $\theta \in S$ such that there exists a minimal prime belonging to the ideal $I_{\theta}$ of $R$ generated by $\{y(\theta, \beta), \beta \in G\}$ is of height at most $1$. Recall that $B$ is the normal closure of $R$.
	
	If not, we can take $S_2=S$ and $n= |B/\mathfrak{m}_B|-1 $ (clearly prime to $p$). Now we fix an element $\theta_1 \in S_2$. For any $\theta' \in S_2$, by Proposition \ref{innovation},  $u(\theta_1, \theta')$ is well-defined and integral over $R$, and hence a unit in $B$. Furthermore, $u(\theta_1, \theta')^n -1$ is an element of the maximal ideal of $B$. Consider the ideal $\mathfrak{b}$ of $B$ generated by $\{u(\theta_1, \theta')^n -1, \theta' \in S_2\}$. By Krull's principal ideal theorem, we can find a prime $\mathfrak{p}$ belonging to $\mathfrak{b}$ of height at most $|S_2|-1$. By Lemma \ref{commutative algebra}, its inverse image $\mathfrak{p}'$ in $R$ is also of height at most $|S_2|-1$. Then we can take $R'=R/\mathfrak{p}'$. The four conditions follow from our construction, Proposition \ref{u} and Proposition \ref{innovation}.
	
	If so, we can choose such an element $\theta \in S$ and consider a prime $\mathfrak{p}_\theta$ (of $R$) of $I_\theta$ of height at most $1$. Take $R_1=R/\mathfrak{p}_\theta$. Then we have $y(\theta, \alpha)=0$ for any $\alpha \in G$ and $\operatorname{dim} R_1 \ge \operatorname{dim} R - 1$.  Now we can use our induction hypothesis for the CNL domain $ R_1$ and the finite set $S\setminus \{\theta\}$. Then we get our desired result.
\end{proof}

\begin{rem}
	This result is also used in \cite{ZHANG2024}. See \cite[Lemma 4.2.5 \& 5.1.2]{ZHANG2024}.
\end{rem}



\section{Galois deformation theory}
In this section, we will study some properties of 2-dimensional universal deformation rings and pseudo-deformation rings.


\subsection{Universal deformation rings}	
In this subsection, we study some properties of universal deformation rings related to $ \rho_x$.\\

Let $F$ be a totally real field. Let $p>3$ be a rational prime and $\mathbb{F}$ be a finite field of characteristic $p$. Let $\Sigma$ be a finite set of finite places of $F$ containing all the places $v \mid p$. Recall that $\mathcal{O}$ is the ring of integers of some local field $E$ (a finite extension of $\mathbb{Q}_p$) with a uniformizer $\pi$ and residue field $\mathbb{F}$. Let $\chi: \operatorname{Gal}(\bar{F}/F) \to \mathcal{O}^\times$ be a continuous totally odd character unramified outside $\Sigma$, and $\bar{\chi}=\chi ~\operatorname{mod} ~\pi$. Let $F_{\Sigma}$ be the maximal extension of $F$ unramified outside $\Sigma$ and all infinite places. We denote $\bar{\rho}_0$ by the representation $1 \oplus \bar{\chi}$.


For an element $0 \ne x \in H^1(F_{\Sigma}/F, \mathbb{F}(\bar{\chi}^{-1}))$, there is a unique representation$$
\rho_x: \operatorname{Gal}(F_{\Sigma}/F) \to \operatorname{GL}_2 (\mathbb{F}), ~~~~~~\rho_x=
\begin{pmatrix}
	1 &  *\\
	~& \bar{\chi}
\end{pmatrix}$$
such that $\rho_x(c)=
\begin{pmatrix}
	1 & ~\\
	~& -1
\end{pmatrix}$ for a complex conjugation $c$, where $*$ corresponds to $x$.

Let $R_x$ be the universal deformation ring in the sense of Mazur \cite{mazur1989deforming} of $ \rho_x$
and $\rho_x^\textnormal{univ}: \operatorname{Gal}(F_{\Sigma}/F) \to \operatorname{GL}_2 (R_x)$ the universal deformation of $ \rho_x$.
Following \cite[Proposition 2.4]{skinner1999residually}, we consider the deformation problem of auxiliary type. 

\begin{defn}\label{auxiliary}
	A deformation $\rho: \operatorname{Gal}(\bar{F}/F) \to \operatorname{GL}_2(A)$ of $ \rho_x$ is of auxiliary type if
	
	1) $A$ is an $\mathcal{O}$-algebra,
	
	2) $\operatorname{det} \rho = \chi$,
	
	3) $\rho$ factors through $\operatorname{Gal}(F_{\Sigma}/F)$.
	\\
	We denote the universal deformation of auxiliary type by $
	\rho_x^\textnormal{aux}: \operatorname{Gal}(F_{\Sigma}/F) \to \operatorname{GL}_2(R_x^\textnormal{aux}).$
\end{defn}

\begin{prop} \label{dimension}  The ring  $R_x^\textnormal{aux}$ is isomorphic to $\mathcal{O}[[x_1,\cdots, x_g]]/(f_1,\cdots, f_r)$, where $g= \operatorname{dim}_{\mathbb{F}} (H^1(F_{\Sigma}/F, \operatorname{ad}^0 \rho_x))$, $ r \le \operatorname{dim}_{\mathbb{F}} (H^2(F_{\Sigma}/F, \operatorname{ad}^0 \rho_x))$.
	
	 
\end{prop}

\begin{proof}
	The assertion is from \cite[Proposition 2]{mazur1989deforming}. 
\end{proof}
 
 \begin{prop}\label{raux}
 	The dimension of every irreducible component of $R_x^\textnormal{aux} $ is at least $2[F: \mathbb{Q}]+1$.
 \end{prop}
 
 \begin{proof}
 	This is a direct consequence of Proposition \ref{dimension}, Krull's principal ideal theorem and the global characteristic formula.
 \end{proof}
 
 Let $\rho_x^\textnormal{aux,red}: \operatorname{Gal}(F_{\Sigma}/F) \to \operatorname{GL}_2(R_x^\textnormal{aux,red})$ be the universal reducible deformation of $\rho_x$ of auxiliary type. Clearly, $R_x^\textnormal{aux,red} $ is a quotient of $R_x^\textnormal{aux} $.
 
 \begin{prop} \label{red} Let $\delta_F$ be the Leopoldt defect (see \cite[10.3.7]{neukirch2013cohomology}. Then we have $\operatorname{dim} (R_x^\textnormal{aux,red}) \le 1+\delta_F +\operatorname{dim}_{\mathbb{F}}(H^1(F_{\Sigma}/F, \mathbb{F}(\bar{\chi}^{-1})))$.
 	
 	%2) $\operatorname{dim} (R_x^\textnormal{aux,red}) \ge 2[F: \mathbb{Q}]+1-\operatorname{dim}_{\mathbb{F}} (H^1(F_{\Sigma}/F, \mathbb{F}(\bar{\chi})))$.
 \end{prop}
 
 \begin{proof}
 	The proof is an analogue to \cite[Section 2.2]{skinner1999residually}. We sketch a proof here. 
 	
 	Let $L(\Sigma)$ is the maximal abelian pro-$p$-extension of $F$ unramified outside $\Sigma$, and write $G=\operatorname{Gal}(L(\Sigma)(\bar{\chi})/F) \simeq \Delta \times \Gamma \times \mathbb{Z}_p^{\delta_F+1}$, where $\Delta = \operatorname{Gal}(F(\bar{\chi})/F))$, $\Gamma$ is a finite $p$-group, $L(\Sigma)(\bar{\chi}) = L(\Sigma)\cdot F(\bar{\chi}) $, and $F(\bar{\chi}) $ is the splitting field. Let $M$ be the maximal abelian pro-$p$-extension of $L(\Sigma)(\bar{\chi})$ unramified outside $\Sigma$ such that $\Delta$ acts on $H=\operatorname{Gal}(M/L(\Sigma)(\bar{\chi}))$ via $\bar{\chi}^{-1}$. Then any reducible deformation of auxiliary type factors through $\operatorname{Gal}(M/F)$.
 	
 	Put $A= \mathbb{Z}_p[[G]]$. By Nakayama's lemma, the group $H$ is a finitely generated $A$-module generated by $m:=\operatorname{dim}_{\mathbb{F}}(H^1(F_{\Sigma}/F, \mathbb{F}(\bar{\chi}^{-1}))) $ elements. Fix a presentation $\mathfrak{a} \to \bigoplus_{i=1}^m Ae_i \twoheadrightarrow H$ such that $e_m$ projects to $h_m \in H $ for which $\rho_x(h_m)=\begin{pmatrix}
 		1 & u\\
 		~ & 1
 	\end{pmatrix}$ with $u \ne 0$ and such that if $i \ne m$ then $e_i$ projects to $h_i \in H $ for which $\rho_x(h_i)=\begin{pmatrix}
 	1 & ~\\
 	~ & 1
 	\end{pmatrix}$. Choose a $u_0 \in \mathcal{O}^\times$ reducing to $u$. Put $A_1 = \mathcal{O}[[\Gamma]][[T_1, \cdots, T_{\delta_F+1}]]$ and fix an embedding $\phi:A \hookrightarrow A \otimes_{\mathbb{Z}_p[\Delta]} \mathcal{O} \simeq A_1$ where the map $\mathbb{Z}_p [\Delta]\to \mathcal{O} $
 	is via $\bar{\chi}$. Let $J$ be the ideal of $A_1[[x_1,\cdots, x_{m-1}]]$ generated by $$
 	\{\phi(a_1)x_1+\cdots+\phi(a_{m-1})x_{m-1}+\phi(a_m)u_0: \Sigma a_ie_i \in \mathfrak{a}\}.$$ 
 	
 	Let $B=A_1[[x_1,\cdots, x_{m-1}]]/J$. Define $\rho: \operatorname{Gal}(M/F) \to \operatorname{GL}_2(B)$ by$$
 	\rho (g)=\begin{pmatrix}
 		\chi(g)\cdot \phi (g)^{-1} & ~\\
 		~ & \phi(g)
 	\end{pmatrix}, ~~~~~~~g \in G, $$
 	$$\rho(h) =\begin{pmatrix}
 	1 & \tau(h)\\
 	~ & 1
 	\end{pmatrix}, \tau(h_i)=e_i, ~ i=1,\cdots,m-1, ~\tau(h_m)=u_0.
 	$$
 	 
 	 By the universal property, we know that $B \cong R_x^\textnormal{aux,red}$. Hence, we have $\operatorname{dim} (R_x^\textnormal{aux,red}) \le 1+\delta_F +1+m-1 =1+\delta_F +\operatorname{dim}_{\mathbb{F}}(H^1(F_{\Sigma}/F, \mathbb{F}(\bar{\chi}^{-1})))$.
 	 
 	 %The second part is a consequence of the results of GMA. We choose a basis for $\rho_x^\textnormal{aux}$ such that $\rho_x^\textnormal{aux}(c)= %\begin{pmatrix}
 	 	%1 & ~\\
 	 	%~& -1
 	 %\end{pmatrix}$ and write $\rho_x^\textnormal{aux}(\sigma)=\begin{pmatrix}
 	 %a_\sigma & b_\sigma\\
 	 %c_\sigma& d_\sigma
 	 %\end{pmatrix}$. Then by Proposition \ref{sub-GMA}, $A= R_x^\textnormal{aux}[\rho_x^\textnormal{aux}( \operatorname{Gal}(F_{\Sigma}/F))]$ is a GMA over $ R_x^\textnormal{aux}$, a sub-$R_x^\textnormal{aux} $-GMA of $M_2(R_x^\textnormal{aux})$.
 	 
 	 %Now write $A= \begin{pmatrix}
 	 %R_x^\textnormal{aux} & B\\
 	 %C & R_x^\textnormal{aux}
 	 %\end{pmatrix}$. Then $C$ is the ideal of $ R_x^\textnormal{aux}$ generated by all $c_\sigma$. By Proposition \ref{GMA} and Nakayama's lemma, $C$ can be generated by at most $\operatorname{dim}_{\mathbb{F}} (H^1(F_{\Sigma}/F, \mathbb{F}(\bar{\chi}))) $ elements as $\operatorname{Ext}^1 ( \bar{\chi}, 1) \cong  H^1(F_{\Sigma}/F, \mathbb{F}(\bar{\chi}))$. By the universal property, we also have $R_x^\textnormal{aux,red} \cong R_x^\textnormal{aux}/C $. Combining Proposition \ref{dimension}, we know that $\operatorname{dim} (R_x^\textnormal{aux,red}) \ge 2[F: \mathbb{Q}]+1-\operatorname{dim}_{\mathbb{F}} (H^1(F_{\Sigma}/F, \mathbb{F}(\bar{\chi})))$.
 	 
 \end{proof}
 
 %\begin{prop}\label{dim aux}
 	%Assume that $F$ is an abelian totally real field. Further assume that $[F:\mathbb{Q}]=\operatorname{dim}_{\mathbb{F}}(H^1(F_{\Sigma}/F, \mathbb{F}(\bar{\chi}^{-1}))) = \operatorname{dim}_{\mathbb{F}} (H^1(F_{\Sigma}/F, \mathbb{F}(\bar{\chi})))$. Then $R_x^\textnormal{aux}$ is a local complete intersection ring of Krull dimension $1+2[F:\mathbb{Q}]$.
 %\end{prop}
 
 %\begin{proof}
 	%From our assumptions and Proposition \ref{red},  the Krull dimension of $R_x^\textnormal{aux,red} $ is $[F:\mathbb{Q}]+1$. By 1) of Proposition \ref{dimension}, we know that $R_x^\textnormal{aux}$ is a local complete intersection ring of Krull dimension $1+2[F:\mathbb{Q}]$.
 %\end{proof}
 
 %\begin{rem}
 	%1) This result is a generalization of \cite[Corollary 3.5]{deo2023density} to abelian totally real fields. In Deo's result, he also assumes the equation of the dimension $1(=[F:\mathbb{Q}])=\operatorname{dim}_{\mathbb{F}}(H^1(\mathbb{Q}_{\Sigma}/\mathbb{Q}, \mathbb{F}(\bar{\chi}^{-1})))$, and in this case, we can get a surjection from the universal pseudo-deformation ring $R^\textnormal{ps} $ to the universal deformation ring $R_x$ by Kisin's argument \cite[Corollary 1.4.4]{kisin2009fontaine} which is a crucial step in Deo's proof. However, we cannot assume that the cohomology group $ H^1(F_{\Sigma}/F, \mathbb{F}(\bar{\chi}^{-1}))$ is of dimension $1$ in the totally real field case because the global characteristic formula tells us that it is of dimension at least $[F:\mathbb{Q}]$. Thus, we have to give a new method to prove it.
 	
 	%2) The condition in our result can be studied by Galois cohomology. See the appendix.
 %\end{rem}
 \begin{prop} \label{char}
 	Assume that for all $v \in \Sigma \setminus \Sigma_p$, we have $p \mid \operatorname{Nm}(v)-1$, and  $\chi$ is unramified outside $ \Sigma_p$.
 	Let $R$ be a CNL domain with maximal ideal $\mathfrak{m}_R$, fraction field $K$ and residue field $\mathbb{F}$. Suppose that $\rho: \operatorname{Gal}(F_{\Sigma}/F) \to \operatorname{GL}_2(R)$ is a deformation of $\rho_x$ of auxiliary type. Then for any $v \in \Sigma \setminus \Sigma_p$, there exists a finite extension $K'/K$ and a finite character $\xi_v: I_v \to K'^{\times}$ of $p$-power order such that $\operatorname{tr}(\rho) \mid_{I_v}= \xi_v+\xi_v^{-1}$.
 \end{prop}
 
 \begin{proof}
 	It is an analogue to \cite[Corollary 5.7.2]{pan2022fontaine}. We sketch a proof here.
 	
 	As $v \nmid p$, we know that for any $\sigma \in I_v$, $\rho(\sigma) $ is of pro-$p$-order. From the structure of Galois group of the maximal tamely ramified extension \cite[Proposition 7.5.2]{neukirch2013cohomology}, we know that $\rho \mid_{I_v}$ only depends on the $\mathbb{Z}_p$-quotient of $I_v $, and we denote the generator by $\sigma_v$.
 	
 	As $\rho$ is of auxiliary type, we just need to show that the eigenvalues $\alpha, \alpha^{-1}$ of $\rho(\sigma_v)$ are of $p$-power order. By \cite[Theorem 7.5.3]{neukirch2013cohomology}, there exists an element in $\tau \in G_{F_v} \setminus I_v$ satisfying $\tau \sigma_v \tau^{-1}=\sigma_v^{\operatorname{Nm}(v)}$. Thus, we have $\alpha^{(\operatorname{Nm}(v)^2-1)}=1$. Since $\alpha \equiv 1 ~\operatorname{mod}~\mathfrak{m}_R$, we know that $\alpha$ is of $p$-power order.
 \end{proof}
 
 Now we suppose that $\mathcal{O}$ has enough $p$-power roots of unity with residue field $\mathbb{F}$. Let $\xi_v: k(v)^\times \to \mathcal{O}^\times$ be characters of $p$-power order for $v \in \Sigma \setminus \Sigma_p$, and we can view them as characters of $I_v$ by local class field theory.
 
 Now consider the deformation $\rho$ of $\rho_x$ such that 
 
 1) it is of auxiliary type,
 
 2) it is unramified outside $\Sigma$,
 
 3) $\operatorname{tr}(\rho) \mid_{I_v}= \xi_v+\xi_v^{-1}$ for all $v \in \Sigma \setminus \Sigma_p$.
 \\
 Let $\rho_x^{\textnormal{aux}, \{\xi_v\}}: \operatorname{Gal}(F_{\Sigma}/F) \to \operatorname{GL}_2(R_x^{\textnormal{aux}, \{\xi_v\}})$ be the universal deformation of such a deformation problem. Obviously, $ R_x^{\textnormal{aux}, \{\xi_v\}}$ is a quotient of $R_x^\textnormal{aux}$. %%For simplicity, we denote $ R_x^{\textnormal{aux}, \{\xi_v\}}$ by $R_x^1$ if $\xi_v$ is trivial for all $v \in \Sigma \setminus \Sigma_p$.
 
 By Proposition \ref{char}, we know that for any minimal prime $\mathfrak{P}$ of $R_x^\textnormal{aux}$, the natural surjection $ R_x^\textnormal{aux} \to R_x^\textnormal{aux}/ \mathfrak{P}$ factors through some $ R_x^{\textnormal{aux}, \{\xi_v\}}$. We say that $ \mathfrak{P}$ \textit{belongs to} $\{\xi_v\}$ in this situation.
 
 
 \subsection{Universal pseudo-deformation ring} \label{mis} In this subsection, we collect some results which will be useful for our pro-modularity arguments.
 
 Keep the notations and assumptions in the previous subsection. Write $\Sigma_p \subset \Sigma$ as the set of places of $F$ lying above $p$. We further assume that for all $v \in \Sigma \setminus \Sigma_p$, we have $p \mid \operatorname{Nm}(v)-1$, and that $\chi: \operatorname{Gal}(F_{\Sigma}/F) \to \mathcal{O}^{\times}$ such that 
 
 1) $\chi$ is unramified outside $ \Sigma_p$,
 
 2) $\chi(\operatorname{Frob}_v) \equiv 1~\operatorname{mod}~\pi$ for  
 $v \in \Sigma \setminus \Sigma_p$.
 
Let $R^\textnormal{ps}$ be the universal pseudo-deformation ring of $\operatorname{tr}(\bar{\rho}_0)= 1+\bar{\chi}$, and $T^\textnormal{univ}$ the universal deformation. The existence of $(R^\textnormal{ps}, T^\textnormal{univ})$ is proved in \cite[Section 2.4]{skinner1999residually}.  Let $ (R_\textnormal{aux}^{\textnormal{ps}}, T_\textnormal{aux}) $ be the universal pseudo-deformation of $\operatorname{tr}(\bar{\rho}_0)$ of auxiliary type (i.e., with $ \operatorname{det} = \chi $)  and $(R^{\textnormal{ps}, \{\xi_v\}}, T^{\textnormal{ps}, \{\xi_v\}})$ the universal pseudo-deformation of $\operatorname{tr}(\bar{\rho}_0)$ of auxiliary type satifying $T_\textnormal{aux}\mid_{I_v}= \xi_v+\xi_v^{-1} $ for all $v \in \Sigma \setminus \Sigma_p$. Clearly, there are natural surjections $R^\textnormal{ps} \twoheadrightarrow R_\textnormal{aux}^\textnormal{ps} \twoheadrightarrow R^{\textnormal{ps}, \{\xi_v\}}$. By \cite[Corollary 5.7.2]{pan2022fontaine}, for any minimal prime $\mathfrak{P}$ of $R_\textnormal{aux}^\textnormal{ps}$, the natural surjection $ R_\textnormal{aux}^\textnormal{ps} \to R_\textnormal{aux}^\textnormal{ps}/ \mathfrak{P}$ factors through some $ R^{\textnormal{ps}, \{\xi_v\}}$. Similarly, we also say that $ \mathfrak{P}$ \textit{belongs to} $\{\xi_v\}$ in this situation.

From our definitions, we can get the following results immediately.

\begin{prop}\label{RX1}
	1) $R_x^{\textnormal{aux}, \{\xi_v\}}/\pi=R_x^{\textnormal{aux}, \{\xi'_v\}}/\pi$.
	
	2) $R^{\textnormal{ps}, \{\xi_v\}}/\pi= R^{\textnormal{ps}, \{\xi'_v\}}/\pi$. 
\end{prop}


Let $(R^{\textnormal{ps,red}}, T^{\textnormal{red}})$ be the universal pseudo-deformation ring of auxiliary type such that $T^{\textnormal{red}}$ is reducible, i.e., $y(\sigma, \tau)=0$ for all $\sigma, \tau \in \operatorname{Gal}(F_{\Sigma}/F)$. 

\begin{prop}\label{red-pseudo} Assume $F$ is abelian over $\mathbb{Q}$.
	
	1) We have $\operatorname{dim}(R^{\textnormal{ps,red}}) \le 2$.
	
	2) Assume $\bar{\chi}$ is quadratic and can be extended to $G_{\mathbb{Q}}$. Then the dihedral locus $C^{\operatorname{dih}}$ of $\operatorname{Spec} R_\textnormal{aux}^{\textnormal{ps}}$, i.e., for any point $\mathfrak{p}$ in $C^{\operatorname{dih}}$, $\rho(\mathfrak{p}) \cong \operatorname{Ind}_{G_L}^{G_F} \theta$ for some character $\theta$ and the splitting field $L=F(\bar{\chi})$, is of dimension at most $2+[F: \mathbb{Q}]$. 
\end{prop}

\begin{proof}
	1) Consider the character $d : \operatorname{Gal}(F_{\Sigma}/F) \to (R^{\textnormal{ps,red}})^\times$ in the sense of Definition \ref{pseudo}, and it induces a natural surjection $\mathcal{O}[[G_{F}^{ab}]] \twoheadrightarrow R^{\textnormal{ps,red}}$. Hence, $\operatorname{dim}(R^{\textnormal{ps,red}}) \le 2+ \delta_F=2$. The last equation is from Leopoldt's conjecture for abelian totally real fields.
	
	2) From our definition, we can find that if $\mathfrak{p} \in C^{\operatorname{dih}}$ and $\rho(\mathfrak{p}) \cong \operatorname{Ind}_{G_L}^{G_F} \theta$, then the global character $\theta$ is a lifting of the character $\bar{\chi}|_{G_L}=\mathbf{1}$. Thus, the dihedral locus is at most of $\operatorname{dim} \mathcal{O}[[G_{L}^{ab}]]$, which is $2+[F: \mathbb{Q}]$ by Leopoldt's conjecture for abelian CM fields.
\end{proof}

\begin{cor}\label{red-one}
	Keep assumptions as in Proposition \ref{red-pseudo}. Let $\mathfrak{s}$ be a prime of $R_\textnormal{aux}^\textnormal{ps} $ containing $p$ and write $R=R_\textnormal{aux}^\textnormal{ps}/\mathfrak{s}$.
	
	1) If $\dim R \ge 2$, then we can find a prime $\mathfrak{s}'$ of $R_\textnormal{aux}^\textnormal{ps} $ containing $\mathfrak{s}$ such that the corresponding representation is irreducible.
	
	2) If $\dim R \ge 2+[F: \mathbb{Q}]$, then we can find a prime $\mathfrak{s}'$ of $R_\textnormal{aux}^\textnormal{ps} $ containing $\mathfrak{s}$ such that the corresponding representation is irreducible and not dihedral.
\end{cor}

\begin{proof}
	We first give a useful lemma in commutative algebra.
	
	\begin{lem}\label{Jacobon}
		Let $R$ be a noetherian local ring with maximal ideal $\mathfrak{m}_R$. Suppose that $V \subsetneqq \operatorname{Spec} R$ is a closed subset. Then there exists a one-dimensional prime $\mathfrak{p} \notin V$.
	\end{lem}
	
	\begin{proof}[Proof of Lemma \ref{Jacobon}]
		Using \cite[Lemma 28.6.4]{stacks-project}, the scheme $\operatorname{Spec} R \setminus\{\mathfrak{m}_R\}$ is Jacobson, i.e., the underlying topological space is Jacobson. As $V \ne \operatorname{Spec} R$, there exists a closed point of $\operatorname{Spec} R \setminus\{\mathfrak{m}_R\}$ not contained in $V$. This shows the result.
	\end{proof}
	
	Our assertions are direct consequences from Lemma \ref{Jacobon} and Proposition \ref{red-pseudo}.
		 %For simplicity, we suppose that $R=R_\textnormal{aux}^\textnormal{ps}/\mathfrak{s}$ is of dimension $2$. For any one-dimensional prime $\mathfrak{t}$ of $R$ such that $T_\textnormal{aux}~ \operatorname{mod}~\mathfrak{t}$ is reducible, we get a surjection $\mathcal{O}[[\operatorname{Gal}(L(\Sigma)/F)]]/(\pi) \twoheadrightarrow R^{\textnormal{ps,red}}/(\pi) \twoheadrightarrow R/\mathfrak{t}$, all of which have dimension $1$ by Proposition \ref{red-pseudo}. Consider an element $t \in \mathfrak{m}_R$ but $t \notin \mathfrak{t}$, and let $\sqrt{(t)}=\cap \mathfrak{t}_i$ be the primary decomposition. If $\mathfrak{t}_i $ is one-dimensional for some $i$, then we can take $\mathfrak{r}=\mathfrak{t}_i $. Otherwise, we get $\sqrt{(t)}=\mathfrak{m}_R$ as $\operatorname{dim} R=2$. Then by Krull's principal ideal theorem, we have $0=\operatorname{dim}(R/tR) \ge \operatorname{dim} R -1 =1$, which is impossible.
\end{proof}

Following \cite{skinner1999residually} and \cite{pan2022fontaine}, it is natural to establish the relationship between deformation rings and pseudo-deformation rings.

\begin{prop}\label{pan}
	1) For any one-dimensional prime ideal $\mathfrak{p}_x \in \operatorname{Spec} R_x$ such that $\rho^\textnormal{univ} $~$\operatorname{mod} $~$\mathfrak{p}_x$ is irreducible, the natural map $\widehat{(R^\textnormal{ps}_{\mathfrak{p}})} \to \widehat{(R_x)_{\mathfrak{p}_x}}$ is an isomorphism, and $\operatorname{Spec}(R_x)_{\mathfrak{p}_x} \to \operatorname{Spec}R^\textnormal{ps}_{\mathfrak{p}}$ is surjective. Here $\mathfrak{p} = \mathfrak{p}_x \cap R^\textnormal{ps}$.
	
	2) For any $\mathcal{Q} \in \operatorname{Spec} R_x $ such that $\rho^\textnormal{univ} $~$\operatorname{mod} $~$\mathcal{Q}$ is irreducible, we have $$
	\operatorname{dim} R_x/\mathcal{Q} \le \operatorname{dim} R^\textnormal{ps}/ (\mathcal{Q} \cap R^\textnormal{ps}).$$ Moreover, if $\mathcal{Q}$ is a minimal prime, so is $\mathcal{Q} \cap R^\textnormal{ps}$.
	
	3) The results in 1), 2) are also true if we replace $(R_x, R^\textnormal{ps} )$ by $ (R_x^\textnormal{aux}, R_\textnormal{aux}^\textnormal{ps})$.
	
\end{prop}

\begin{proof}
	The first two assertions are \cite[Corollary 2.3.6]{pan2022fontaine}. The proof of the last assertion is the same as Pan's.
	
\end{proof}

\begin{defn} \cite[Definition 5.4.3]{pan2022fontaine}
	Let $R$ be a CNL ring. The connectedness dimension of $R$ is defined to be 
	$$c(R):=\min_{Z_1,Z_2}\{\dim (Z_1\cap Z_2)\}, $$
	where $Z_1,Z_2$ are non-empty unions of irreducible components of $\operatorname{Spec} R$ such that $Z_1\cup Z_2=\operatorname{Spec} R$.
\end{defn}

\begin{prop}\label{connect dim} Suppose $ \mathfrak{p}$ is a maximal ideal of $R_\textnormal{aux}^\textnormal{ps}[1/p]$ whose corresponding representation is absolutely irreducible.
	
	1) $\operatorname{dim}  (R_\textnormal{aux}^\textnormal{ps})_\mathfrak{p} \ge 2[F: \mathbb{Q}]$.
	
	2) The connectedness dimension $c((R_\textnormal{aux}^\textnormal{ps})_\mathfrak{p})$ is at least $2[F: \mathbb{Q}]-1$.
\end{prop}

\begin{proof}
	For the first part, see \cite[Lemma 7.1.2]{pan2022fontaine}.
	For the second part, see \cite[Lemma 7.4.6]{pan2022fontaine}.
\end{proof}

\begin{rem}
	Unlike the universal deformation ring $R_x^\textnormal{aux} $, the structure of the universal pseudo-deformation ring $R_\textnormal{aux}^\textnormal{ps} $ is mysterious because we cannot use the global characteristic formula as in Proposition \ref{dimension}. However, we can show that if $\mathfrak{p}$ is a one-dimensional irreducible prime of $R_\textnormal{aux}^\textnormal{ps} $, then every irreducible component of $\operatorname{Spec} R_\textnormal{aux}^\textnormal{ps}$ containing $\mathfrak{p}$ is of dimension at least $ 2[F: \mathbb{Q}]+1$, as an analogue to Proposition \ref{raux}. The proof is as follows.
	
	If our claim is not true, then we have a partition $Z_1\amalg Z_2=\operatorname{Spec} (R_\textnormal{aux}^\textnormal{ps})_\mathfrak{p}$, where $Z_i$ are finite sets of some irreducible components of $ \operatorname{Spec} (R_\textnormal{aux}^\textnormal{ps})_\mathfrak{p}$ such that $Z_1$ consists of all ones of dimension at least $ 2[F: \mathbb{Q}]$. By the first part of Proposition \ref{connect dim} and our assumption, both $Z_1$ and $Z_2$ are non-empty. By the second part of Proposition \ref{connect dim}, we can find $C_1 \in Z_1, C_2 \in Z_2$ and a prime $\mathfrak{q} \in C_1 \cap C_2$ of dimension at least $2[F: \mathbb{Q}]-1$. As $\operatorname{dim} C_2 \le 2[F: \mathbb{Q}]-1$, we know that $\mathfrak{q}$ is the generic point of $C_2$, which is against to the fact $\mathfrak{q} \in C_1$. Thus, $Z_2$ has to be empty.
	
	We do not know whether every irreducible component of $R_\textnormal{aux}^\textnormal{ps} $ is of dimension at least $ 2[F: \mathbb{Q}]+1$.
\end{rem}

\section{Pro-modularity}
In this section, we will prove the pro-modularity of the universal deformation ring. We fix a totally real field $F$ and a prime $p > 3$.

\subsection{Hecke algebras}
In this subsection, we introduce Hecke operators and Hecke algebras. The main reference is \cite[Section 3]{pan2022fontaine}.

Assume that the prime $p$ splits completely in the totally real field $F$.

Let $D$ be a quaternion algebra with centre $F$ which is ramified at all infinite places of $F$ and unramified at all places above $p$, and we fix isomorphisms $D \otimes F_v \cong M_2(F_v)$ for any $v$ where $D$ is unramified. We view $K_p = \prod_{v \mid p} \operatorname{GL}_2(\mathcal{O}_{F_v})$ and $D_p^\times =  \prod_{v \mid p} \operatorname{GL}_2(F_v)$ as subgroups of $(D \otimes_F \mathbb{A}_F)^\times$.

Let $A$ be a topological $\mathcal{O}$-algebra and $U = \prod_{v \mid p} U_v$ be an open compact subgroup of $(D \otimes_F \mathbb{A}_F^\infty)^\times$ such that $U_v \subseteq \operatorname{GL}_2(\mathcal{O}_{F_v})$ for $v \mid p$. We write $U^p= \prod_{v \nmid p} U_v$ (tame level) and $U_p =\prod_{v \mid p} U_v$. Let $\psi: (\mathbb{A}_F^\infty)^\times/F_+^\times \to A^\times$ be a continuous character, where $F_+$ is the set of totally positive elements in $F$. Let $\tau: \prod_{v \mid p} U_v \to \operatorname{Aut}(W_\tau)$ be a continuous representation on a finite $A$-module $W_\tau$, and we also view $\tau$ as a representation of $U$ by projecting to $\prod_{v \mid p} U_v$. 

Suppose that $\xi: U^p \to A^\times$ is a continuous smooth character such that $\psi \mid_{\prod_{v \nmid p}(\mathcal{O}_{F_v}^\times) \cap U_v}= \xi \mid_{\prod_{v \nmid p}(\mathcal{O}_{F_v}^\times) \cap U_v}$. We define $S_{\tau, \psi, \xi}(U,A)$ as the space of continuous functions: $f: D^\times \setminus (D \otimes_F \mathbb{A}_F^\infty)^\times \to W_\tau$, such that for any $g \in (D \otimes_F \mathbb{A}_F^\infty)^\times$, $z \in (\mathbb{A}_F^\infty)^\times$, $u=u^pu_p \in U$, we have $$
f(guz)= \psi(z)\xi(u^p)\tau(u_p^{-1})(f(g)).$$ 
If $\psi \mid_{U_p \cap \mathcal{O}_{F, p}^\times} = \tau^{-1} \mid_{U_p \cap \mathcal{O}_{F, p}^\times}$, then we have $$
S_{\tau, \psi, \xi}(U,A) \simeq \bigoplus_{i \in I} W_\tau ^{(t_i^{-1}D^{\times}t_i \cap U (\mathbb{A}_F^\infty)^\times)/F^\times},$$
where $I=D^\times \setminus (D \otimes_F \mathbb{A}_F^\infty)^\times/ U(\mathbb{A}_F^\infty)^\times$ and $\{t_i\}_{i \in I}$ is a set of representatives. We will simply write $S_{\psi, \xi}(U,A)$ if $\tau$ is the trivial action on $A$. We say that $U$ is \textit{sufficiently small} if $ (t_i^{-1}D^{\times}t_i \cap U (\mathbb{A}_F^\infty)^\times)/F^\times$ is trivial for all $i \in I$.

Let $\Sigma$ be a finite set of primes of $F$ containing all places above $p$ and places $v$ where either $D$ is ramified or $U_v$ is not a maximal open subgroup. For any $v \notin \Sigma$, we define the Hecke operator $T_v \in \operatorname{End}(S_{\tau, \psi, \xi}(U,A))$ to be the double coset action $[U_v \begin{pmatrix}
	\pi_v & ~\\
	~ & 1
\end{pmatrix}U_v]$. Precisely, if we write $U_v \begin{pmatrix}
\pi_v & ~\\
~ & 1
\end{pmatrix}U_v  = \coprod_i \gamma_i U_v $, then $$
(T_v \cdot f) (g) = \sum_i f(g\gamma_i), ~~f \in S_{\tau, \psi, \xi}(U,A). $$

We define Hecke algebra $\mathbb{T}_{\tau, \psi, \xi} ^\Sigma (U,A) \subseteq \operatorname{End}(S_{\tau, \psi, \xi}(U,A))$ to be the $A$-subalgebra generated  by all $T_v, v \notin \Sigma$. This is a finite commutative $A$-algebra. We will write $\mathbb{T}_{\psi, \xi}^\Sigma (U,A)$ for the Hecke algebra if $\tau $ is the trivial action on $A$. Suppose that $A= \mathcal{O}/ \pi^n$, $ U$ is sufficiently small and $\psi \mid_{U \cap (\mathbb{A}_F^\infty)^\times} $ is trivial modulo $ \pi^n$. Then we can check that the Hecke algebra $\mathbb{T}_{\psi, \xi}^\Sigma (U, \mathcal{O}/ \pi^n)$ is independent of $\Sigma$, and we may drop $\Sigma$ in this situation. See \cite[3.3.3]{pan2022fontaine}.

\begin{defn} (Big Hecke algebra) \label{big}
	Let $U^p$ be a tame level and $\psi: (\mathbb{A}_F^\infty)^\times/ (U^p \cap (\mathbb{A}_F^\infty)^\times)F_+^\times \to \mathcal{O}^\times$ be a continuous character. Then we define the \textit{Hecke algebra} $$
	\mathbb{T}_{\psi, \xi}(U^p) = \varprojlim_{(n, U_p) \in \mathcal{I}} \mathbb{T}_{\psi, \xi}(U^pU_p, \mathcal{O}/\pi^n), $$
where $\mathcal{I}$ is the set of pairs $(n, U_p)$ with $U_p \subseteq K_p$ and $n$ a positive integer such that $ \psi \mid_{U_p \cap \mathcal{O}_{F,p}^\times} \equiv 1~ \operatorname{mod}~\pi^n$.
	
\end{defn}


\subsection{Deformation rings and Hecke algebras}\label{sec 2}
In this subsection, we collect some results about deformation rings and Hecke algebras, most of which are from \cite{pan2022fontaine}.

Assume $F$ is an abelian totally real number field of even degree over $\mathbb{Q}$ in which $p$ splits completely. As in Section \ref{mis}, we assume that for all $v \in \Sigma \setminus \Sigma_p$, we have $p \mid \operatorname{Nm}(v)-1$. Let $\xi_v: k(v)^\times \to \mathcal{O}^\times$ be characters of $p$-power order for $v \in \Sigma \setminus \Sigma_p$, and we can view them as characters of $I_v$ by local class field theory. Let $\chi: \operatorname{Gal}(F_{\Sigma}/F) \to \mathcal{O}^{\times}$ be a continuous totally odd character such that 

1) $\chi$ is unramified outside $ \Sigma_p$,

2) $\chi(\operatorname{Frob}_v) \equiv 1~\operatorname{mod}~\pi$ for  
$v \in \Sigma \setminus \Sigma_p$.
\\

Let $D$ be a quaternion algebra over $F$ ramified exactly at all infinite places, and we fix an isomorphism between $( D \otimes_F \mathbb{A}_F^\infty)^\times$ and $\operatorname{GL}_2(\mathbb{A}_F^\infty)$. Let $\psi=\chi \epsilon$ be a character of $ (\mathbb{A}_F^\infty)^\times/ F_+^\times$ via global class field theory. Let $U^p=\prod_{v \nmid p} U_v$ be a tame level such that $U_v= \operatorname{GL}_2(\mathcal{O}_{F_v})$ if $v \notin \Sigma$ and $U_v = \operatorname{Iw}_v := \{g \equiv \begin{pmatrix}
	* & *\\
	0 & *
\end{pmatrix}
~\operatorname{mod}~\pi_v, g \in \operatorname{GL}_2(\mathcal{O}_{F_v})\}$ otherwise. For any $v \in \Sigma \setminus \Sigma_p$, the map $\begin{pmatrix}
	a & b\\
	c & d
\end{pmatrix} \to \xi_v (a/d~\operatorname{mod}~\pi_v)$ defines a character of $U_v$ and the product of $\xi_v$ can be viewed as a character $\xi$ of $U^p$ by projecting to $\prod_{v \in  \Sigma \setminus \Sigma_p} U_v$. From our setting, we can define a Hecke algebra $\mathbb{T}_{\xi} := \mathbb{T}_{\psi, \xi}(U^p)$ as in the previous subsection.\\

From now on, we suppose that $\bar{\rho}_0= 1 \oplus \bar{\chi}$ is modular, i.e. $T_v-(1+\chi(\operatorname{Frob}_v)), v \notin \Sigma$ and $\pi$ generate a maximal ideal $\mathfrak{m}_{\xi}$ of $ \mathbb{T}_{\xi}$. 

\begin{thm} \label{dimofT}
	Each irreducible component of $(\mathbb{T}_{\xi})_{\mathfrak{m}_{\xi}} $ is of characteristic zero and of dimension at least $1+2[F:\mathbb{Q}]$.
\end{thm}

\begin{proof}
	See \cite[Theorem 3.6.1]{pan2022fontaine}.
\end{proof}
	
	\begin{rem}
		In some cases, the big Hecke algebra is equidimensional of dimension $1+2[F:\mathbb{Q}]$. See the proof of \cite[Theorem 3.6.1]{pan2022fontaine} and \cite[Remark 3.6.10]{pan2022fontaine}.
	\end{rem}
	
As the discussion in \cite[3.3.4]{pan2022fontaine}, there is a natural pseudo-deformation $T_{\mathfrak{m}_{\xi}} : \operatorname{Gal}(F_{\Sigma}/F) \to (\mathbb{T}_{\xi})_{\mathfrak{m}_{\xi}}$ with determinant $\chi$ sending $\operatorname{Frob}_v$ to Hecke operator $T_v$ for $v \notin \Sigma$. By the universal property, we get a natural surjection $ R_\textnormal{aux}^\textnormal{ps} \twoheadrightarrow (\mathbb{T}_{\xi})_{\mathfrak{m}_{\xi}}$. By the local-global compatibility at $v \in \Sigma \setminus \Sigma_p$, we know that such surjection factors through $R^{\textnormal{ps}, \{\xi_v\}} $. Hence, we also get a natural surjection $R^{\textnormal{ps}, \{\xi_v\}} \twoheadrightarrow (\mathbb{T}_{\xi})_{\mathfrak{m}_{\xi}} $.

For any $\mathfrak{p} \in \operatorname{Spec} (\mathbb{T}_{\xi})_{\mathfrak{m}_{\xi}} $, we can define a two-dimensional semi-simple representation $\rho(\mathfrak{p}): \operatorname{Gal}(F_{\Sigma}/F) \to \operatorname{GL}_2(k(\mathfrak{p}))$ with trace $T_{\mathfrak{m}_{\xi}}~\operatorname{mod}~\mathfrak{p}$ in the sense of Definition \ref{construction}.

The following definitions are from \cite[Section 4.1]{pan2022fontaine}.

\begin{defn}\label{nice}
	(1) We say that a prime  of $ R^{\textnormal{ps}, \{\xi_v\}}$ is \textit{pro-modular} if it comes from a prime of $(\mathbb{T}_{\xi})_{\mathfrak{m}_{\xi}} $. We say that a prime of $ R_\textnormal{aux}^\textnormal{ps}$ is \textit{pro-modular} if it comes from a pro-modular prime of some $ R^{\textnormal{ps}, \{\xi_v\}}$.
	
	(2) Let $\mathfrak{q} $ be a prime of $ (\mathbb{T}_{\xi})_{\mathfrak{m}_{\xi}} $ and $A$ be the normal closure of $(\mathbb{T}_{\xi})_{\mathfrak{m}_{\xi}} / \mathfrak{q} $ in $k(\mathfrak{q})$. We say that $\mathfrak{q} $ is a \textit{nice} prime if $\mathfrak{q} $ contains $p$ and $\operatorname{dim} (\mathbb{T}_{\xi})_{\mathfrak{m}_{\xi}} / \mathfrak{q}=1$ and there exists a two-dimensional representation $\rho(\mathfrak{q})^o : \operatorname{Gal}(F_{\Sigma}/F) \to \operatorname{GL}_2(A)$ satisfying:
	
	\quad 1) $\rho(\mathfrak{q})^o \otimes k(\mathfrak{q}) \cong \rho(\mathfrak{q})$ is irreducible.
	
	\quad 2) The $\operatorname{mod}~\mathfrak{m}_A$ reduction $\bar{\rho}_b$ of $\rho(\mathfrak{q})^o$ is a non-split extension and has the form $ \bar{\rho}_b(g)=\begin{pmatrix}
		* & *\\
		0 & *
	\end{pmatrix}$, $g \in \operatorname{Gal}(F_{\Sigma}/F)$. Here $\mathfrak{m}_A $ is the maximal ideal of $A$.
	
	\quad 3) (dihedral condition) If $\rho(\mathfrak{q}) $ is dihedral, namely isomorphic to $\operatorname{Ind}_{G_L}^{G_F} \theta$ for some quadratic extension $L$ of $F$ and continuous character $\theta: G_L \to k(\mathfrak{q})^\times$, then $L \cap F(\zeta_p) = F$,
	where $\zeta_p $ is a primitive $p$-th root of unity.
	
	\quad 4) $\rho(\mathfrak{q})^o \mid_{G_{F_v}} = \bar{\rho}_b \mid_{G_{F_v}} $ for any $v \in \Sigma \setminus \Sigma_p$.
	
	(3) We say that a prime of $ R^{\textnormal{ps}, \{\xi_v\}}$ is \textit{nice} if it comes from a nice prime of $(\mathbb{T}_{\xi})_{\mathfrak{m}_{\xi}} $.
	
\end{defn}

The following proposition gives some sufficient conditions for the third condition of (2) in Definition \ref{nice}.

\begin{prop}\label{equiv}
	Let $\mathfrak{q}$ be a prime ideal of $ R^{\textnormal{ps}, \{\xi_v\}}$ containing $p$ such that $ R^{\textnormal{ps}, \{\xi_v\}}/\mathfrak{q}$ is one-dimensional. Suppose $\rho(\mathfrak{q})$ is irreducible. Then the third condition of (2) in Definition \ref{nice} holds for $\mathfrak{q}$ if one of the following conditions holds:
	
	1)  $\bar{\chi}$ is not quadratic. 
	
	2) $\bar{\chi}\mid_{G_{F_v}}$ is trivial for some $v \mid p$.
	
	3) There exists a place $v \mid p$ such that $\bar{\chi}\mid_{G_{F_v}}$ is not trivial and $\rho(\mathfrak{q})^o\mid_{G_{F_v}} \cong \begin{pmatrix}
		\chi_{v,1} & *\\
		0 & \chi_{v,2}
	\end{pmatrix}$ is reducible. Moreover, $ \chi_{v,1}$ is of infinite order.
\end{prop}

\begin{proof}
	See \cite[Lemma 4.1.6]{pan2022fontaine}.
\end{proof}

The following theorem is a key evidence to prove our  pro-modularity argument.

\begin{thm} \label{R=T}   $(R_{\mathfrak{q}}=\mathbb{T}_{\mathfrak{q}})$
	Let $\mathfrak{q} $ be a nice prime of $(\mathbb{T}_{\xi})_{\mathfrak{m}_{\xi}} $ and $\mathfrak{q}^\textnormal{ps}= \mathfrak{q} \cap R^{\textnormal{ps}, \{\xi_v\}}$. Then the natural surjective map $ (R^{\textnormal{ps}, \{\xi_v\}})_{\mathfrak{q}^\textnormal{ps}} \twoheadrightarrow (\mathbb{T}_{\xi})_{\mathfrak{q}}$ has nilpotent kernel.
	
\end{thm}

\begin{proof}
	See \cite[Section 4]{pan2022fontaine}.
\end{proof}

\begin{rem}\label{dim of nice}
	We briefly summarize some main steps of the proof of Theorem \ref{R=T}. Here we follow most notations as in \cite[Section 4]{pan2022fontaine}, and one can find more details there.
	
	Let $\mathfrak{q}$ be a nice prime of $(\mathbb{T}_{\xi})_{\mathfrak{m}_{\xi}}$. Let $B$ be the topological closure of the $\mathbb{F}$-algebra generated by all the entries of $\rho(\mathfrak{q})^o(\operatorname{Gal}(F_{\Sigma}/F))$. By Chebotarev’s density Theorem, we may find a finite set of primes $T'$ disjoint with $\Sigma$ such
	that the entries of $\rho(\mathfrak{q})^o(\operatorname{Frob}_v), v \in T'$, topologically generate $B$. Let $P= T' \cup \Sigma$.
	
	Let $r=\operatorname{dim}_{k(\mathfrak{q})} H^1(G_{F,P},\operatorname{ad}^0 \rho(\mathfrak{q})(1))$ (the cardinality of Taylor-Wiles primes). Let $R_{\infty}^{\{\xi_v\}}= R_{\textnormal{loc}}^{\{\xi_v\}}[[x_1, ..., x_g]]$, where $ R_{\textnormal{loc}}^{\{\xi_v\}}$ is the completed tensor product of local deformation rings determined by primes in $P$ (see \cite[Definition 4.2.2]{pan2022fontaine}) and $g=r+|P|-[F:\mathbb{Q}]-1$. Let $S_\infty'=\mathcal{O}[[y_1, ... , y_{4|P|-1}, s_1', ... , s_r']]$, and $S_\infty=\mathcal{O}[[y_1, ... , y_{4|P|-1}, s_1, ... , s_r]]$ with an ideal $\mathfrak{a}_1=(y_1, ... , y_{4|P|-1}, s_1, ... , s_r)$ such that $ S_\infty$ is a finite free $S_\infty'$-algebra.
	
	Let $\textnormal{m}_\infty^{\{\xi_v\}}$ and $\textnormal{m}_0^{\{\xi_v\}}$ be the \textit{patched Hecke modules }(see \cite[Section 4.7]{pan2022fontaine}). Then $\textnormal{m}_\infty^{\{\xi_v\}} $ is a flat  $ S_\infty$-module and $\textnormal{m}_\infty^{\{\xi_v\}}/ \mathfrak{a}_1 \textnormal{m}_\infty^{\{\xi_v\}} \cong (\textnormal{m}_0^{\{\xi_v\}})^{\oplus 2^r}$. By Pan's local-global compatibility result \cite[Corollary 3.5.8]{pan2022fontaine}, $\textnormal{m}_0^{\{\xi_v\}}$ is a finitely generated faithful $\widehat{(\mathbb{T}_{\xi})_{\mathfrak{q}}}$-module. By Theorem \ref{dimofT}, we have $\operatorname{dim}_{\widehat{(\mathbb{T}_{\xi})_{\mathfrak{q}}}} (\textnormal{m}_0^{\{\xi_v\}})= \operatorname{dim} \widehat{(\mathbb{T}_{\xi})_{\mathfrak{q}}} \ge 2[F:\mathbb{Q}]$ since $\mathfrak{q}$ is of dimension $1$.
	
	Let $(R^{\{\xi_v\}})'$ be $\widehat{(R_{\infty}^{\{\xi_v\}})_{\mathfrak{q}_{\infty}^{\{\xi_v\}}}} \otimes_{S_\infty'} S_\infty$, where ${\mathfrak{q}}_{\infty}^{\{\xi_v\}} $ is the prime of $ R_{\infty}^{\{\xi_v\}}$ corresponding to the nice prime $\mathfrak{q}$.	We know that $\widehat{(\mathbb{T}_{\xi})_{\mathfrak{q}}} $ is a natural quotient of $(R^{\{\xi_v\}})'$ by mapping $ S_\infty$ to $S_\infty/ \mathfrak{a}_1=\mathcal{O}$. Note that $ y_1, ... , y_{4|P|-1}, s_1, ... , s_r \in \mathfrak{a}_1$ form a regular sequence of $\textnormal{m}_\infty^{\{\xi_v\}}$. Then we can conclude that $$\operatorname{dim}_{(R^{\{\xi_v\}})'}(\textnormal{m}_\infty^{\{\xi_v\}}) \ge 4|P|-1 +r +\operatorname{dim} \widehat{(\mathbb{T}_{\xi})_{\mathfrak{q}}} \ge 4|P|-1 +r+2[F:\mathbb{Q}].$$
	
	Using Taylor's Ihara avoidance trick as in \cite{PMIHES_2008__108__183_0}, one can show that  
	$\widehat{(R_{\infty}^{\{\xi_v\}})_{\mathfrak{q}_{\infty}^{\{\xi_v\}}}}$ is equidimensional of dimension $ 4|P|-1 +r+2[F:\mathbb{Q}]$ (see \cite[Section 4.2, Lemma 4.8.2]{pan2022fontaine}). As $(R^{\{\xi_v\}})'$ is finite free over $\widehat{(R_{\infty}^{\{\xi_v\}})_{\mathfrak{q}_{\infty}^{\{\xi_v\}}}}$, we have $$\operatorname{dim}_{(R^{\{\xi_v\}})'}(\textnormal{m}_\infty^{\{\xi_v\}})=\operatorname{dim}_{\widehat{(R_{\infty}^{\{\xi_v\}})_{\mathfrak{q}_{\infty}^{\{\xi_v\}}}}}(\textnormal{m}_\infty^{\{\xi_v\}})=4|P|-1 +r+2[F:\mathbb{Q}]= \operatorname{dim} \widehat{(R_{\infty}^{\{\xi_v\}})_{\mathfrak{q}_{\infty}^{\{\xi_v\}}}}.$$
	
	If the characters in $\{\xi_v\}$ are all non-trivial, by \cite[Proposition 4.2.3]{pan2022fontaine},  $\widehat{(R_{\infty}^{\{\xi_v\}})_{\mathfrak{q}_{\infty}^{\{\xi_v\}}}}$ is irreducible. From the previous equations, we know that $\textnormal{m}_\infty^{\{\xi_v\}} $ has full support on $\widehat{(R_{\infty}^{\{\xi_v\}})_{\mathfrak{q}_{\infty}^{\{\xi_v\}}}}$. By \cite[Lemma 4.8.4]{pan2022fontaine}, Theorem \ref{R=T} holds in this case. Similar to the proof of \cite[Theorem 4.1]{PMIHES_2008__108__183_0}, one can deduce the general case using Taylor's trick.
	
	From Theorem \ref{dimofT}, \cite[Corollary 12.5]{eisenbud2013commutative}, and discussions above, we can also conclude that for a nice prime $\mathfrak{q}$ of $(\mathbb{T}_{\xi})_{\mathfrak{m}_{\xi}}$, we have $\operatorname{dim} \widehat{(\mathbb{T}_{\xi})_{\mathfrak{q}}} =\operatorname{dim} (\mathbb{T}_{\xi})_{\mathfrak{q}}= 2[F:\mathbb{Q}]$. In other words, if an irreducible component of the big Hecke algebra contains a nice prime, then it is of dimension $1+2[F:\mathbb{Q}]$.
\end{rem}

\begin{rem}\label{remark}
	From the definitions and the "$R=\mathbb{T}$" theorem above, we can conclude the following two facts.
	
	1) If a minimal prime $\mathfrak{P} \in \operatorname{Spec} R^{\textnormal{ps}, \{\xi_v\}}$ is pro-modular, then any prime contained in the irreducible component defined by $\mathfrak{P}$ is pro-modular. We also say that an irreducible component of $\operatorname{Spec} R^{\textnormal{ps}, \{\xi_v\}}$ is pro-modular if its generic point is pro-modular.
	
	2) If an irreducible component of $\operatorname{Spec} R^{\textnormal{ps}, \{\xi_v\}}$ contains a nice prime, then it is pro-modular.
\end{rem}

\subsection{The ordinary case}\label{sec 3}
In this subsection, we recall some results on the deformation rings and Hecke algebras in the ordinary case.

Assume that $F$ is an abelian totally real number field over $\mathbb{Q}$ in which $p$ is unramified. Let $\Sigma$ be a finite set of primes of $F$ containing all places above $p$. Let $\chi: \operatorname{Gal}(F_{\Sigma}/F) \to \mathcal{O}^{\times}$ be a continuous totally odd character such that 

1) $\bar{\chi}$ can be extended to a character of $G_{\mathbb{Q}}$.

2) $\bar{\chi} \mid_{G_{F_v}}$ is not trivial for any $v \mid p$.

3) $\chi \mid_{G_{F_v}}$ is de Rham for any $v \mid p$.
\\

Let $(R^{\textnormal{ps,ord}}, T^{\textnormal{ord}})$ be the universal pseudo-deformation ring of $\operatorname{tr}(\bar{\rho}_0)$ of auxiliary type such that $ T^{\textnormal{ord}}\mid_{G_{F_v}}$ is reducible for any $v \mid p$, i.e. $y(\sigma, \tau)=0$ for any $\sigma, \tau \in G_{F_v}$, $v \mid p$.

Since $\bar{\chi} \mid_{G_{F_v}}$ is not trivial, we have $T^{\textnormal{ord}}\mid_{G_{F_v}}= \psi_{v,1}^{\textnormal{ord}}+\psi_{v,2}^{\textnormal{ord}}$ for some characters $ \psi_{v,1}^{\textnormal{ord}}, \psi_{v,2}^{\textnormal{ord}}: G_{F_v} \to (R^{\textnormal{ps,ord}})^\times$ which are liftings of $ \textbf{1}, \bar{\chi}\mid_{G_{F_v}}$ respectively, and actually, $\psi_{v,1}^{\textnormal{ord}} = a\mid_{G_{F_v}} $ and $\psi_{v,2}^{\textnormal{ord}} = d\mid_{G_{F_v}} $ in the sense of Definition \ref{pseudo}. Hence, $ \psi_{v,1}^{\textnormal{ord}}\mid_{I_v}$ induces a homomorphism $\mathcal{O}[[\mathcal{O}_{F_v}^\times(p)]] \to R^{\textnormal{ps,ord}}$ for any $v \mid p$ via local class field theory, where $\mathcal{O}_{F_v}^\times(p) $ is the $p$-adic completion of $ \mathcal{O}_{F_v}^\times$.

\begin{defn}\cite[Definition 5.3.1]{pan2022fontaine}
	Let $T: \operatorname{Gal}(F_\Sigma/F) \to R$ be a two-dimensional pseudo-representation over some ring $R$ such that for some place $v$, $T \mid_{G_{F_v}} = \psi_1 + \psi_2$ is a sum of two characters. We say $T$ is $\psi_1$-\textit{ordinary} if for any $\sigma, \tau \in G_{F_v}$, $\eta \in \operatorname{Gal}(F_\Sigma/F)$, $$
	T(\sigma \tau \eta)-\psi_1(\sigma)T(\tau \eta)-\psi_2(\tau)T(\sigma\eta)+\psi_1(\sigma)\psi_2(\tau)T(\eta)=0.$$
\end{defn}

\begin{prop}\label{psi1}
	Let $\rho : \operatorname{Gal}(F_\Sigma/F) \to  \operatorname{GL}_2(R)$ be a two-dimensional irreducible representation over some ring $R$ with trace $T$ such that $T \mid_{G_{F_v}} = \psi_1 + \psi_2$, a sum of two characters. Suppose $\rho\mid_{G_{F_v}}$ has the form $\begin{pmatrix}
		\psi_{1} & *\\
		0 & \psi_{2}
	\end{pmatrix}$. Then $T$ is $\psi_1$-ordinary. Conversely, if $R$ is a field and $T$ is $\psi_1$-ordinary, then after possibly enlarging $R$,
	we have $\rho\mid_{G_{F_v}} \cong \begin{pmatrix}
		\psi_{1} & *\\
		0 & \psi_{2}
	\end{pmatrix}$.
\end{prop}

\begin{proof}
	See \cite[Lemma 5.3.2]{pan2022fontaine}.
\end{proof}

  Define the \textit{Iwasawa algebra} $$
\Lambda_F := \widehat{\bigotimes}_{v \mid p} \mathcal{O}[[\mathcal{O}_{F_v}^\times(p)]], $$
and we get a natural map $\Lambda_F \to R^{\textnormal{ps,ord}} $ induced by the characters $ \psi_{v,1}^{\textnormal{ord}}$. The following result is the ordinary Fontaine-Mazur conjecture proved in \cite{skinner1999residually} and \cite[Section 5]{pan2022fontaine}, which is used to produce pro-modular primes in our setting.

\begin{prop}\label{ordinary}
	(1) $ R^{\textnormal{ps,ord}}$ is a finite $ \Lambda_F$-algebra. 
	
	(2) For any maximal ideal $\mathfrak{p}$ of $R^{\textnormal{ps,ord}}[\frac{1}{p}]$, we denote the associated semi-simple representation $\operatorname{Gal}(F_{\Sigma}/F) \to \operatorname{GL}_2(k(\mathfrak{p}))$ by $\rho(\mathfrak{p})$. Assume that
	
	\quad i) $ \rho(\mathfrak{p})$ is irreducible.
	
	\quad ii) For any $v \mid p$, $\rho(\mathfrak{p}) \mid_{G_{F_v}} \cong  \begin{pmatrix}
		\chi_{v,1} & *\\
		0 & \chi_{v,2}
	\end{pmatrix}$ such that $\chi_{v,1} $ is de Rham and has strictly less Hodge-Tate number than $\chi_{v,2} $ for any embedding $F_v \hookrightarrow \overline{\mathbb{Q}}_p$.
	
	Then $ \rho(\mathfrak{p})$ comes from a twist of Hilbert modular form.
\end{prop}

\begin{proof}
	See \cite[Theorem 5.1.1]{pan2022fontaine}.
\end{proof}

\begin{lem}\label{dense}
	Let $C^{\textnormal{ord}}$ be an irreducible component of $ R^{\textnormal{ps,ord}}$ of dimension at least $1+[F: \mathbb{Q}]$.
	 Let $ C^{\textnormal{ord,aut}}$ be the set of regular de Rham primes in $C^{\textnormal{ord}}$. Then $C^{\textnormal{ord,aut}} $ is dense in $C^{\textnormal{ord}}$. Here a prime $\mathfrak{q}$ is called regular de Rham prime if
	
	\quad i) $p \notin \mathfrak{q}$ and $R^{\textnormal{ps,ord}}/\mathfrak{q} $ is one-dimensional.
	
	\quad ii) The semi-simple representation $\rho(\mathfrak{q}): \operatorname{Gal}(F_{\Sigma}/F) \to \operatorname{GL}_2(k(\mathfrak{q}))$ (in the sense of Definition \ref{pseudo}) is irreducible.
	
	\quad iii) For any $v \mid p$, $\rho(\mathfrak{q}) \mid_{G_{F_v}} \cong  \begin{pmatrix}
		\chi_{v,1} & *\\
		0 & \chi_{v,2}
	\end{pmatrix}$ such that $\chi_{v,1} $ is de Rham and has strictly less Hodge-Tate number than $\chi_{v,2} $ for any embedding $F_v \hookrightarrow \overline{\mathbb{Q}}_p$.
\end{lem}

\begin{proof}
	The proof is the same as \cite[Corollary 7.2.3]{pan2022fontaine}. We sketch it here.
	
	From our assumption and Proposition \ref{ordinary}, we know that the morphism $C^{\textnormal{ord}} \to \operatorname{Spec} \Lambda_F$. 
	
	  For $v \mid p$, let $ \psi_{v,1}^{\textnormal{ord}}, \psi_{v,2}^{\textnormal{ord}}: G_{F_v} \to (R^{\textnormal{ps,ord}})^\times$ be the liftings of $1, \bar{\chi}\mid_{G_{F_v}}$ respectively. By Proposition \ref{psi1}, there exists $n_v \in \{1,2\}$ such that for any $\mathfrak{q} \in C_1^{\textnormal{ord}}$, we get $\rho(\mathfrak{q}) \mid_{G_{F_v}} \cong  \begin{pmatrix}
		\psi_{v,n_v}^{\textnormal{ord}} ~\operatorname{mod}~\mathfrak{q}& *\\
		0 & \psi_{v,3-n_v}^{\textnormal{ord}}~\operatorname{mod}~\mathfrak{q}
	\end{pmatrix}$. 
	
	Note that the map $\Lambda_F \to R^{\textnormal{ps,ord}} $ is induced by the characters $ \psi_{v,1}^{\textnormal{ord}}$. Hence, we know that the image of $ C^{\textnormal{ord,aut}}$ is dense in $\operatorname{Spec} \Lambda_F$, and hence dense in $C^{\textnormal{ord}}$.
\end{proof}


%\begin{cor}\label{dimord}
	 %The Krull dimension of $ R^{\textnormal{ps,ord}}$ is $1+[F:\mathbb{Q}]$.
%\end{cor}

%\begin{proof}
	%This result is based on \cite[Section 5]{pan2022fontaine}. We just sketch a proof here.
	
	%By Proposition \ref{ordinary}, we know that $\operatorname{dim} R^{\textnormal{ps,ord}} \le 1+[F:\mathbb{Q}]$. If we consider the ordinary Hecke algebra $\mathbb{T}^{\textnormal{ord}}$(see \cite[Section 5.2]{pan2022fontaine} for the explicit definition), by \cite[Corollary 5.2.1]{pan2022fontaine}, it is equidimensional of dimension $1+[F:\mathbb{Q}]$. As there is a two-dimensional pseudo-representation $T^{\textnormal{ord}}: \operatorname{Gal}(F_\Sigma/F) \to \mathbb{T}^{\textnormal{ord}}$, we get a natural surjection $ R^{\textnormal{ps,ord}} \twoheadrightarrow \mathbb{T}^{\textnormal{ord}}$, and we have $\operatorname{dim} R^{\textnormal{ps,ord}} \ge 1+[F:\mathbb{Q}]$. Hence, $\operatorname{dim} R^{\textnormal{ps,ord}} = 1+[F:\mathbb{Q}]$.
%\end{proof}



\subsection{The pro-modularity of universal deformation rings}
In  this subsection, we will prove our pro-modularity statements.

Keep the notations and assumptions on the finite set $\Sigma$ and the character $\chi$ in the previous two subsections. In this subsection, we further assume that $\bar{\chi}\mid_{G_{F_v}} \ne \omega_p^{\pm 1}$ for any $v\mid p$, and $[F:\mathbb{Q}] \ge 7|\Sigma \setminus \Sigma_p|+2.$

\begin{defn}\label{pro-modular}
	We say that a prime of $ R_x^{\textnormal{aux}, \{\xi_v\}}$ is pro-modular if its inverse image of the natural map $ R^{\textnormal{ps}, \{\xi_v\}} \to R_x^{\textnormal{aux}, \{\xi_v\}}$ is pro-modular. We say that a prime of $ R_x^\textnormal{aux}$ is pro-modular if it is the inverse image of some pro-modular prime of some $ R_x^{\textnormal{aux}, \{\xi_v\}}$ via the natural surjection $ R_x^\textnormal{aux} \twoheadrightarrow R_x^{\textnormal{aux}, \{\xi_v\}}$. 
\end{defn}

\begin{thm}\label{pro}
	For any irreducible component of $R_\textnormal{aux}^\textnormal{ps} $ of dimension at least $1+2[F:\mathbb{Q}]$, it is pro-modular.
\end{thm}

\begin{proof}
	Let $C$ be an irreducible component of $R_\textnormal{aux}^\textnormal{ps} $ of dimension at least $1+2[F:\mathbb{Q}]$. Consider the corresponding generic point $\mathfrak{P}$ of $R_\textnormal{aux}^\textnormal{ps} $, and we suppose that it belongs to $\{\xi_v\}$. We claim that $\mathfrak{P}$ is pro-modular.
	
	For convenience for readers, we briefly introduce our strategy here. By Remark \ref{remark}, we just need to find a nice prime containing $\mathfrak{P}$. In other words, we need to find a pro-modular prime which satisfies the conditions we list in (2) of Definition \ref{nice}. Hence, our proof can be divided into two parts. In the first part, following \cite[Section 7]{pan2022fontaine}, we use Pan's results given in the previous section to find enough pro-modular primes from ordinary points. (By the way, there is also a similar step in the proof of Deo's big $R=\mathbb{T}$ theorem. See \cite[Proposition 3.4]{deo2023density}.) In the second part, we use Corollary \ref{innovation'} to find a nice prime, which is a main innovation of this paper.\\
	
    Write $C^{\textnormal{ord}}=\operatorname{Spec} R^{\textnormal{ps,ord}} \cap C$.
	
	\begin{lem}\label{cord}
		$\operatorname{dim} C^{\textnormal{ord}} \ge [F:\mathbb{Q}]+1$.
	\end{lem} 
	
	\begin{proof}
		For any $v \mid p$, let $R_v^\textnormal{ps}$ be the universal pseudo-deformation ring of the pseudo-representation of $G_{F_v} $ lifting $1+\bar{\chi}\mid_{G_{F_v}}$ of auxiliary type, i.e. of determinant $\chi \mid_{G_{F_v}}$. Let $R_v^{\textnormal{ps,ord}}$ be the quotient of $R_v^\textnormal{ps}$ parametrizing all reducible liftings. As $\bar{\chi} \ne 1, \omega_p^{\pm 1}$, we know that the kernel of $ R_v^\textnormal{ps} \twoheadrightarrow R_v^{\textnormal{ps,ord}}$ is a principal ideal by \cite[Corollary B.20]{pavskunas2013image}. Write $R_p^\textnormal{ps}=\widehat{\bigotimes}_{\mathcal{O}} R_v^\textnormal{ps}$ and $R_p^{\textnormal{ps,ord}}=\widehat{\bigotimes}_{\mathcal{O}} R_v^{\textnormal{ps,ord}}$. Then by universal property, we get $R^{\textnormal{ps,ord}} =R_\textnormal{aux}^\textnormal{ps} \otimes_{R_p^\textnormal{ps}} R_p^{\textnormal{ps,ord}}$. Thus, $\operatorname{dim} C^{\textnormal{ord}} \ge \operatorname{dim} C -  |\Sigma_p | \ge 1+ 2[F:\mathbb{Q}]-[F:\mathbb{Q}]=[F:\mathbb{Q}]+1$.
	\end{proof}
	
	
	
	Now we can choose an irreducible component $C_1^{\textnormal{ord}} \subset C^{\textnormal{ord}}$ of $\operatorname{Spec} R^{\textnormal{ps,ord}}$ of dimension $1+[F: \mathbb{Q}]$, and we let $\mathfrak{Q}$ be the generic point of $C_1^{\textnormal{ord}}$. For any $v \in \Sigma \setminus \Sigma_p$, we fix a lifting $\sigma_v $ of the generator of the $\mathbb{Z}_p$-quotient of $I_v$ and a lifting $\operatorname{Frob}_v$ of the Frobenius element. Write $S=\{\sigma_v, \operatorname{Frob}_v:~v \in \Sigma \setminus \Sigma_p\}$. Clearly, $|S|=2|\Sigma \setminus \Sigma_p|$.
	
	Let $I$ be an ideal of $R_\textnormal{aux}^\textnormal{ps}/(\mathfrak{Q}, \pi)$ generated by $$\{a(\sigma_{v})-1, a(\operatorname{Frob}_{v})-1, d(\operatorname{Frob}_{v})-1,~ v\ \in \Sigma \setminus \Sigma_p\}.$$ 
	We consider a minimal prime $\mathfrak{u}$ of $I$. By Krull's principal ideal theorem, $(R_\textnormal{aux}^\textnormal{ps}/(\mathfrak{Q}, \pi))/\mathfrak{u}$ is of dimension at least $[F: \mathbb{Q}]-3|\Sigma \setminus \Sigma_p|$. For simplicity, we write such a domain as $R_0$.\\
	
	To find a nice prime, we consider the following steps.\\
	
	\textit{Step} I. We start from the CNL domain $R_0$.
	
	By Corollary \ref{innovation'}, we can find a partition of $S = S_1 \amalg S_2$, a positive integer $n>1$ prime to $p$ and a CNL domain $R_0'$ satisfying the following conditions.
	
	1) $R_0'$ is a quotient of $R_0$.
	
	2) For any $\theta \in S_1$, we have $y(\theta, \alpha)=0$ for any $\alpha \in \operatorname{Gal}(F_\Sigma/F)$ in $R_0'$. For any $\theta', \theta'' \in S_2$, we have $y(\theta', \alpha)^n=y(\theta'', \alpha)^n$ for any $\alpha \in \operatorname{Gal}(F_\Sigma/F)$ in $R_0'$.
	
	3) For any $\alpha \in \operatorname{Gal}(F_\Sigma/F), \theta \in S_2$, either $ y(\theta, \beta)=0$ for any $\beta \in \operatorname{Gal}(F_\Sigma/F)$ or $u(\alpha, \theta)$ is well-defined and integral over $R_0'$.
	
	4) We have $\operatorname{dim} R_0' \ge \operatorname{dim} R_0 - |S|$. If $S_2$ is not empty, then further $\operatorname{dim} R_0' \ge \operatorname{dim} R_0 - |S|+1$.\\
	
	\textit{Step} II. We start from the CNL domain $R_0'$.
	
	Similar to \textit{Step} I, we can find a partition of $S = S_1' \amalg S_2'$, a positive integer $n'>1$ prime to $p$ and a CNL domain $R_0''$ satisfying the following conditions.
	
	1') $R_0''$ is a quotient of $R_0'$.
	
	2') For any $\theta \in S_1'$, we have $y(\alpha, \theta)=0$ for any $\alpha \in \operatorname{Gal}(F_\Sigma/F)$ in $R_0''$. For any $\theta', \theta'' \in S_2'$, we have $y(\alpha, \theta')^{n'}=y(\alpha, \theta'')^{n'}$ for any $\alpha \in \operatorname{Gal}(F_\Sigma/F)$ in $R_0''$.
	
	3') For any $\alpha \in \operatorname{Gal}(F_\Sigma/F), \theta \in S_2'$, either $ y(\beta, \theta)=0$ for any $\beta \in \operatorname{Gal}(F_\Sigma/F)$ or $u'(\alpha, \theta)$ is well-defined and integral over $R_0''$.
	
	4') We have $\operatorname{dim} R_0'' \ge \operatorname{dim} R_0' - |S|$. If $S_2'$ is not empty, then further $\operatorname{dim} R_0'' \ge \operatorname{dim} R_0' - |S|+1$.\\
	 
	 \textit{Step} III. We consider the domain $R_0''$.
	 
	 If both $S_2$ and $S_2'$ are not empty, then we fix an element $\theta_{i} \in S_2$ and an element $\theta_{j} \in S_2'$. Consider a prime $\mathfrak{p}_{ij}$ of $R_0''$ of the ideal $(y(\theta_{i}, \theta_{j}))$ and let $R''=R_0''/\mathfrak{p}_{ij}$. By Proposition \ref{u}, we know that either $y(\theta_{i}, \alpha)=0$ or $y(\beta, \theta_{j})=0$ for all $\alpha, \beta \in \operatorname{Gal}(F_\Sigma/F)$. By Krull's principal ideal theorem, we have $\operatorname{dim} R'' \ge \operatorname{dim} R_0'' -1$.
	 
	 If not, we take $R''=R_0''$.\\
	 
	 In conclusion, one of the following three assertions holds for $R''$.
	 
	 a) $ y(\theta, \alpha) =y(\alpha, \theta)=0$ for all $\theta \in S$ and $\alpha \in \operatorname{Gal}(F_\Sigma/F)$.
	 
	 b) There exists an element $\theta_{i} \in S$ such that either $y(\theta, \alpha)=0$ or $y(\theta, \alpha)^{n}=y(\theta_i, \alpha)^{n}$ for all $\theta \in S$ and $\alpha \in \operatorname{Gal}(F_\Sigma/F)$, and $y(\alpha, \theta) =0 $ for all $\theta \in S$ and $\alpha \in \operatorname{Gal}(F_\Sigma/F)$.
	 
	 c) There exists an element $\theta_{j} \in S$ such that either $y(\alpha, \theta)=0$ or $y(\alpha, \theta)^{n'}=y(\alpha, \theta_j)^{n'}$ for all $\theta \in S$ and $\alpha \in \operatorname{Gal}(F_\Sigma/F)$, and $y(\theta, \alpha) = 0 $ for all $\theta \in S$ and $\alpha \in \operatorname{Gal}(F_\Sigma/F)$.
	 
	 In particular, we have $\operatorname{dim} R'' \ge [F: \mathbb{Q}]-7|\Sigma \setminus \Sigma_p| \ge 2$.
	 
	\begin{lem}
		There exists a nice prime containing $\mathfrak{P}$.
	\end{lem}
	
	\begin{proof}
		By Proposition \ref{ordinary}, for any prime $\mathfrak{q}_1 \in  C_1^{\textnormal{ord,aut}}$, the subset of $C_1^{\textnormal{ord}} $ consisting of all regular de Rham primes as defined in Lemma \ref{dense}, we know that $\rho(\mathfrak{q}_1)$ comes from a twist of Hilbert modular form. By the density result in Lemma \ref{dense}, the image of any prime of $C_1^{\textnormal{ord}}$ in $\operatorname{Spec} R^{\textnormal{ps},\{\xi_v\}}$ is pro-modular.  
		
		By Corollary \ref{red-one}, we can choose an irreducible one-dimensional prime $\mathfrak{r}'$ of $R''$. We have shown that $\mathfrak{r}' $ is pro-modular. Now we claim that $\mathfrak{r}' $ is actually a nice prime in the sense of Definition \ref{nice}.
		
		Let $A$ be the normal closure of $R''/\mathfrak{r}'$. Then we obtain a continuous irreducible representation $\rho(\mathfrak{r}') : \operatorname{Gal}(F_{\Sigma}/F) \to \operatorname{GL}_2(A)$ associated to the pseudo-deformation in the following form $$
		\rho(\mathfrak{r}')(\sigma)=\begin{pmatrix}
			a(\sigma) & \frac{y(\sigma, \tau_{\mathfrak{r}'})}{y(\sigma_{\mathfrak{r}'}, \tau_{\mathfrak{r}'})}\\
			y(\sigma_{\mathfrak{r}'}, \sigma)& d(\sigma)
	\end{pmatrix}, ~y(\sigma_{\mathfrak{r}'}, \tau_{\mathfrak{r}'}) \ne 0.$$
		
		Now we verify the conditions in Definition \ref{nice} for all the three cases. 
		
		In case a), we can take $\rho(\mathfrak{r}')^o=\rho(\mathfrak{r}')$. From our construction, the first three conditions are clear (using 3) of Proposition \ref{equiv}). By Proposition \ref{RX1}, we know that if $a(\sigma_{v})=1$, then we have $d(\sigma_{v})=1$ for any $v \in \Sigma \setminus \Sigma_p$ as $T(\sigma_{v})=\xi_v+\xi_v^{-1} \equiv 2~ \operatorname{mod}~\pi$. Thus, we can find that for all $v \in \Sigma \setminus \Sigma_p$, $\rho(\mathfrak{r}')^o \mid_{G_{F_v}}= \begin{pmatrix}
			1 & 0\\
			0 & 1
		\end{pmatrix}$. Then $\mathfrak{r}' $ is a nice prime.
		
		In case b), if $y(\theta_{i}, \alpha)=0$ for all $\alpha \in \operatorname{Gal}(F_\Sigma/F)$, then it is the same as case a). If not, by Proposition \ref{u}, we know that $y(\theta_{i}, \tau_{\mathfrak{r}'}) \ne  0$, otherwise $y(\sigma_{\mathfrak{r}'}, \tau_{\mathfrak{r}'}) = 0$. By 3) in \textit{Step} I, $u(\sigma_{\mathfrak{r}'}, \theta_{i})$ is well-defined and integral over $R''/\mathfrak{r}'$, and hence a unit in $A$. Then we can take $\rho(\mathfrak{r}')^o=P\rho(\mathfrak{r}')P^{-1}$, where $$P= \begin{pmatrix}
			u(\sigma_{\mathfrak{r}'}, \theta_{i}) & 0\\
			0 & 1
		\end{pmatrix} \in \operatorname{GL}_2(A).$$
		
		Similar to the case a), we only need to check the fourth condition. From our construction, we know that either $u(\theta, \sigma_{\mathfrak{r}'})=0$ or $u(\theta, \theta_i)^{n}=1$ for any $\theta \in S$. As $(n,p)=1$, by Hensel's lemma, $ u(\theta, \theta_{i})$ is an element of $\mathbb{F}' \hookrightarrow \mathbb{F}'[[T]]=A$, where $\mathbb{F}' $ is the residue field of $A$. Thus, the fourth condition holds.
		
		In case c), we consider the continuous irreducible Galois representation $\rho(\mathfrak{r}')' : \operatorname{Gal}(F_{\Sigma}/F) \to \operatorname{GL}_2(A)$ in the following form $$
		\rho(\mathfrak{r}')'(\sigma)=\begin{pmatrix}
			a(\sigma) & y(\sigma, \tau_{\mathfrak{r}'})\\
			\frac{y(\sigma_{\mathfrak{r}'}, \sigma)}{y(\sigma_{\mathfrak{r}'}, \tau_{\mathfrak{r}'})}& d(\sigma)
		\end{pmatrix}, ~y(\sigma_{\mathfrak{r}'}, \tau_{\mathfrak{r}'}) \ne 0.$$
		
		If $y(\alpha, \theta_j)=0$ for all $\alpha \in \operatorname{Gal}(F_\Sigma/F)$, then we can take $$
		\rho(\mathfrak{r}')^o=\begin{pmatrix}
			0 & 1\\
			1 & 0
		\end{pmatrix} \rho(\mathfrak{r}')'\begin{pmatrix}
		0 & 1\\
		1 & 0
		\end{pmatrix} \subset \operatorname{GL}_2(A). $$
		Similar to case a), we can vefiry that $\mathfrak{r}'$ is a nice prime.
		If not, we can take $$ \rho(\mathfrak{r}')^o=\begin{pmatrix}
			0 & u'(\tau_{\mathfrak{r}'}, \theta_{j})\\
			1 & 0
		\end{pmatrix} \rho(\mathfrak{r}')'\begin{pmatrix}
			0 & 1\\
			u'(\theta_{j}, \tau_{\mathfrak{r}'}) & 0
		\end{pmatrix} \subset \operatorname{GL}_2(A). $$
	Similar to case b), we can also vefiry that $\mathfrak{r}'$ is a nice prime.
	\end{proof}
	
	From the previous lemma and Theorem \ref{R=T}, we know that $C$ is pro-modular.
	
\end{proof}

\begin{rem}\label{reason}
	1)  By Proposition \ref{connect dim}, for a one-dimensional prime $\mathfrak{p}$ of $R_\textnormal{aux}^\textnormal{ps} $, if it corresponds to an irreducible pseudo-representation, then at least one of the irreducible components of $R_\textnormal{aux}^\textnormal{ps} $ containing $\mathfrak{p}$ is of dimension at least $1+2[F:\mathbb{Q}]$. Hence $\mathfrak{p}$ is pro-modular.
	
	2) In the proof of Theorem \ref{pro}, we can also find that the irreducible component $C$ is of dimension $1+2[F:\mathbb{Q}]$. Note that the Iwasawa algebra $\Lambda_F$ is of dimension $1+[F:\mathbb{Q}]$. By the finiteness result (Proposition \ref{ordinary}), we have $\operatorname{dim} C^{\textnormal{ord}} \le [F:\mathbb{Q}]+1$. Combining Lemma \ref{cord}, we have $\operatorname{dim} C^{\textnormal{ord}} = [F:\mathbb{Q}]+1$, and hence $ \operatorname{dim} C = 2[F:\mathbb{Q}]+1$.
	
	%(2) In the case $\bar{\chi}=\omega_p^{\pm1}$, we can also prove a pro-modularity result as Theorem \ref{pro}, but the deduction is more complicated. Our method of finding nice primes seems to be helpful to prove the two-dimensional Fontaine-Mazur conjecture in the missing case of \cite[Theorem 1.0.2]{pan2022fontaine}. More precisely, to find nice primes, we need to start from a pro-modular prime of large dimension. In Pan's setting, the dimension $[F: \mathbb{Q}]$ is large enough as we can ignore the ``dihedral condition" by Lemma \ref{equiv} ($\omega_p^{-1}$ is not quadratic for $p>3$). However, as $\omega_p$ is quadratic when $p=3$, we need to ensure that the Galois representation associated to the (nice) prime that we construct is not dihedral. Hence, the dimension $[F: \mathbb{Q}]$ seems to be not large enough because the upper bound of the dimension of the modulo $p$ dihedral locus of the universal pseudo-deformation ring is $[F: \mathbb{Q}]+1$ by Leopoldt's conjecture in the abelian CM case. Nevertheless, our strategy here of finding nice primes (the utilization of Corollary \ref{innovation'}) is effective to overcome this difficulty so that we just need to find one-dimensional irreducible pro-modular primes instead of nice primes (see \cite[Lemma 5.1.2]{ZHANG2024}). Thus, the dimension $[F: \mathbb{Q}]$ is still large enough in the missing case. See \cite{ZHANG2024} for more details.
\end{rem}


As a consequence of our pro-modularity argument, we can get a conditional $R=\mathbb{T}$ theorem.

\begin{thm}\label{MINIMAL}
	 If we further assume $$[F:\mathbb{Q}] \ge \operatorname{max}\{1+\frac{1}{2}\operatorname{dim}_{\mathbb{F}}(H^1(F_{\Sigma}/F, \mathbb{F}(\bar{\chi}^{-1}))), 7|\Sigma \setminus \Sigma_p|+2\},$$ 
	then for any non-zero element $x \in H^1(F_{\Sigma}/F, \mathbb{F}(\bar{\chi}^{-1}))$, any prime ideal of $ R_x^\textnormal{aux}$ is pro-modular. Moreover, $ R_x^\textnormal{aux}$ is a local complete intersection ring of Krull dimension $1+2[F:\mathbb{Q}] $.
	
	%%$(2)$ Keep the assumption as in $(1)$. If the natural surjection $R^{\textnormal{ps}, \{\xi_v\}} \twoheadrightarrow (\mathbb{T}_{\xi_v})_{\mathfrak{m}_{\xi_v}}$ factors through $R_x^{\textnormal{aux}, \{\xi_v\}} $ for some non-zero cocycle $x \in H^1(F_{\Sigma}/F, \mathbb{F}(\bar{\chi}^{-1}))$, then we have $$(R_x^{\textnormal{aux}, \{\xi_v\}})^{\operatorname{red}}  \cong ((\mathbb{T}_{\xi_v})_{\mathfrak{m}_{\xi_v}})^{\operatorname{red}}$$.
\end{thm}

\begin{proof}
	Let $\mathfrak{P}$ be a minimal prime of $ R_x^\textnormal{aux}$. Let $\mathfrak{P}'$ be the inverse image of $\mathfrak{P}$ via the natural map $f: R_\textnormal{aux}^\textnormal{ps} \to R_x^\textnormal{aux}$. From our assumptions, Proposition \ref{raux}, Proposition \ref{red} and Proposition \ref{pan}, we know that $\mathfrak{P}'$ is a minimal prime of dimension at least $1+2[F:\mathbb{Q}] $. By Theorem \ref{pro}, we know that $\mathfrak{P}$ is pro-modular. 
	
	By Remark \ref{reason}, we know that $ \mathfrak{P}'$ is of dimension $1+2[F:\mathbb{Q}] $. Consequently, $ R_x^\textnormal{aux}$ is of dimension at most $1+2[F:\mathbb{Q}] $. Combining Proposition \ref{dimension} and Proposition \ref{raux}, we obtain that $R_x^\textnormal{aux}$ is a local complete intersection ring of Krull dimension $1+2[F:\mathbb{Q}]$. This proves our result.
	
	%%If the map $R^{\textnormal{ps}, \{\xi_v\}} \twoheadrightarrow (\mathbb{T}_{\xi_v})_{\mathfrak{m}_{\xi_v}}$ factors through $R_x^{\textnormal{aux}, \{\xi_v\}} $, then we know that the map $r: R_x^{\textnormal{aux}, \{\xi_v\}} \to (\mathbb{T}_{\xi_v})_{\mathfrak{m}_{\xi_v}}$ is also a surjection since $(\mathbb{T}_{\xi_v})_{\mathfrak{m}_{\xi_v}} $ is topologically generated by Hecke operators $T_v$, $v \notin \Sigma$. Consider a minimal prime $\mathfrak{p} \in \operatorname{Spec} R_x^{\textnormal{aux}, \{\xi_v\}} $. By (1), we know that $\mathfrak{p}$ actually comes from a prime $ \mathfrak{p}' \in \operatorname{Spec}  (\mathbb{T}_{\xi_v})_{\mathfrak{m}_{\xi_v}}$. That is to say, we have $ R_x^{\textnormal{aux}, \{\xi_v\}}/r^{-1}(\mathfrak{p}') \cong (\mathbb{T}_{\xi_v})_{\mathfrak{m}_{\xi_v}}/\mathfrak{p}' \hookrightarrow R_x^{\textnormal{aux}, \{\xi_v\}}/\mathfrak{p}$. As $\mathfrak{p}$ is minimal, we have $R_x^{\textnormal{aux}, \{\xi_v\}}/\mathfrak{p} \cong (\mathbb{T}_{\xi_v})_{\mathfrak{m}_{\xi_v}}/\mathfrak{p}'$. Thus, $\operatorname{ker} r$ is contained in any minimal prime of $R^{\textnormal{ps}, \{\xi_v\}}$, and this proves (2).
\end{proof}

\subsection{A special case}
In this subsection, we deal with the case $\bar{\chi}\mid_{G_{F_v}} = \omega_p$ for any $v|p$. In this case, the deduction is more complicated because Lemma \ref{cord} no longer holds. We follow the strategy of the proof in \cite[Section 7.4]{pan2022fontaine}.

Keep assumptions on as in Section \ref{sec 2} and Section \ref{sec 3} the finite set $\Sigma$ and the character $\chi$ in the previous two subsections. We further assume $[F:\mathbb{Q}] \ge 7|\Sigma \setminus \Sigma_p|+4.$

\begin{prop}\label{promodular to nice}
	Let $\mathfrak{r}$ be a one-dimensional irreducible pro-modular prime of $R_\textnormal{aux}^\textnormal{ps} $. Then every irreducible component of $R_\textnormal{aux}^\textnormal{ps} $ containing $\mathfrak{r}$ contains a nice prime, and hence is pro-modular.
\end{prop}

\begin{proof}
	Let $Z$ be the set of all irreducible components of $(R_\textnormal{aux}^\textnormal{ps})_{\mathfrak{r}} $. Write $Z= Z_1 \amalg Z_2$, where $Z_1$ is the subset of $Z$ consisting of all pro-modular irreducible components. We first show that if $Z_1$ is not empty, then $Z_2$ is empty.
	
	If not, by Proposition \ref{connect dim}, there exist an irreducible component $Y \in Z_2$ and a pro-modular prime $\mathfrak{r}' \in Y$ of dimension at least $2[F: \mathbb{Q}]-1$.
	
	Arguing as the proof of Theorem \ref{pro}, we may find a prime $\mathfrak{r}'' \subset \mathfrak{r}'$ satisfying the following conditions:
	\begin{itemize}
		\item $\varpi \in \mathfrak{r}''$.
		\item  The dimension of $\mathfrak{r}'' $ is at least $ 2[F: \mathbb{Q}]-7|\Sigma\setminus\Sigma_p|-2 \ge [F: \mathbb{Q}]+2$.
		\item  For any irreducible one-dimensional prime $\mathfrak{t}$ containing $\mathfrak{r}''$, there exists a $2$-dimensional Galois representation $\rho(\mathfrak{t}) : G_{F, \Sigma} \to \operatorname{GL}_2(k(\mathfrak{t}))$ satisfying the conditions in (2) of Definition \ref{nice} except the third one.
	\end{itemize}
	
	If $\bar{\chi}$ is not quadratic, then by Lemma \ref{equiv}, we can find a nice prime containing $\mathfrak{r}''$ (hence containing $\mathfrak{r}'$). Using Theorem \ref{R=T}, we conclude that $Y$ is pro-modular, which contradicts to our definition of $Z_2$.
	
	If $\bar{\chi}$ is quadratic, we know that the splitting field $F(\bar{\chi})$ is an abelian CM field. Then by Corollary \ref{red-one}, we may find a prime containing $\mathfrak{r}'$, whose corresponding Galois representation is irreducible and not dihedral. From our choice of $\mathfrak{r}''$, we know that this is actually a nice prime. Using Theorem \ref{R=T} again, we get a contradiction.
	
	As $\mathfrak{r}$ is pro-modular, using Corollary \ref{dimofT}, we know that there exists a pro-modular prime $\mathfrak{r}_1$ contained in $\mathfrak{r}$ of dimension at least $1+2[F: \mathbb{Q}]$. Arguing as above, we can find a nice prime containing $\mathfrak{r}_1$. This implies that  at least one of the irreducible component containing $\mathfrak{r}$ is pro-modular. Hence, $Z_1$ is not empty.
\end{proof}

From now on, we assume $\bar{\chi}\mid_{G_{F_v}} = \omega_p$ for any $v|p$.
Suppose that $C$ is an irreducible component of $R_\textnormal{aux}^\textnormal{ps} $ of dimension at least $1+2[F:\mathbb{Q}]$. In this case, by \cite[Lemma 7.4.4]{pan2022fontaine}, the kernel of the natural surjection $ R_v^\textnormal{ps} \twoheadrightarrow R_v^{\textnormal{ps,ord}}$ can be generated by two elements. Then by Krull's principal ideal theorem, we have $\operatorname{dim} C^{\textnormal{ord}} \ge 1$. 

Choose a prime $\mathfrak{q} \in C^{\textnormal{ord}}$ such that $\dim R_\textnormal{aux}^\textnormal{ps}/\mathfrak{q}=1$. Enlarging $\mathcal{O}$ if necessary, we may assume the normalization of $R_\textnormal{aux}^\textnormal{ps}/\mathfrak{q}$ is either $\mathcal{O}$ or isomorphic to $\mathbb{F}[[T]]$. There are three possibilities for the associated semi-simple representation $\rho(\mathfrak{q})$:
\begin{enumerate}
	\item $\rho(\mathfrak{q})$ is irreducible.
	\item $\rho(\mathfrak{q})\cong\psi_1\oplus \psi_2$ is reducible  and $\psi_1/\psi_2$ is not of the form $\epsilon^{\pm 1}\theta$, where $\theta$ is a finite order character of $G_{F}$. We call this case \textit{generic reducible}.
	\item $\rho(\mathfrak{q})\cong\psi_1\oplus \psi_2$ and $\psi_1/\psi_2=\varepsilon\theta$, for some finite order character $\theta$ of $G_{F}$. We call this case \textit{non-generic reducible}.
\end{enumerate}

\begin{thm}\label{special case}
	The irreducible component $C$ is pro-modular of dimension $1+2[F:\mathbb{Q}]$.
\end{thm}

\begin{proof}
	As the discussions above, we divide the proof into three cases.
	
	\textit{1. The irreducible case}
	
	Using the same proof of \cite[Lemma 7.4.7]{pan2022fontaine}, we find that the localization $(R^{\textnormal{ps,ord}})_{\mathfrak{q}}$ is of dimension at least $[F:\mathbb{Q}]$. In other words, there exists an irreducible component of $ R^{\textnormal{ps,ord}}$ containing $\mathfrak{q}$ of dimension at least $1+[F:\mathbb{Q}]$. Arguing as in the proof of Theorem \ref{pro}, we know that $\mathfrak{q}$ is pro-modular. By Proposition \ref{promodular to nice} and Remark \ref{dim of nice}, we know that $C$ is pro-modular of dimension $1+2[F:\mathbb{Q}]$.
	
	\textit{2. The generic reducible case}
	
	By the same proof as in \cite[Section 7.4.4]{pan2022fontaine}, we have $\operatorname{dim} C^{\textnormal{ord}} \ge 1+[F:\mathbb{Q}]$ (see \cite[Corollary 7.4.9]{pan2022fontaine}). Thus, as in the irreducible case, we also have $C$ is pro-modular of dimension $1+2[F:\mathbb{Q}]$.
	
	\textit{3. The non-generic reducible case}
	
	We only sketch the proof. 
	
	For any non-zero class $B\in \operatorname{Ext}^1_{E[G_{F,\Sigma}]}(\psi_1,\psi_2) \cong H^1(G_{F,\Sigma},E(\psi_2/\psi_1))$, it corresponds to a non-split extension $\rho_B: G_{F,\Sigma}\to\operatorname{GL}_2(E)$ of $\psi_1$ by $\psi_2$. We may assume $\rho_B$ is of the form
	\[\begin{pmatrix} \psi_2 & b_B \\ 0 & \psi_1\end{pmatrix}.\] We denote the universal deformation ring of $\rho_B$ with fixed determinant $\chi$ by $R_B$. 
	
	We define $Z_B$ as the set of irreducible components of $R_\textnormal{aux}^\textnormal{ps}$ whose generic points lie in the image of $\operatorname{Spec} R_{B}\to \operatorname{Spec} R_\textnormal{aux}^\textnormal{ps}$. 
	
	First, by the arguments in \cite[Step (1), page 1153]{pan2022fontaine}, there exists a non-zero element $B_0$ in $H^1(G_{F,\Sigma},E(\psi_2/\psi_1)) $ such that we can find a pro-modular irreducible component in $Z_{B_0}$. Using \cite[Corollary 7.4.13]{pan2022fontaine} and Proposition \ref{promodular to nice}, we know that any irreducible component in  $Z_{B_0}$ is pro-modular.
	
	Second, using the same proof in \cite[Step (2), page 1154 \& 1155]{pan2022fontaine}, we can find an element $B_1 $ in $H^1(G_{F,\Sigma},E(\psi_2/\psi_1)) $ such that $C \in Z_{B_1}$ and	there exist irreducible components $C_{B_0}\in Z_{B_0}, C_{B_1}\in Z_{B_1}$ containing a prime of dimension $\mathfrak{s}$ at least $2+[F: \mathbb{Q}]$. As $C_{B_0} $ is pro-modular, we know that $\mathfrak{s}$ is pro-modular. By Corollary \ref{red-one}, we can find a one-dimensional irreducible pro-modular prime in $C_{B_0} \cap C_{B_1} $. Combining Proposition \ref{promodular to nice}, we know that $C_{B_1} $ is pro-modular. Using \cite[Corollary 7.4.13]{pan2022fontaine} and Proposition \ref{promodular to nice} again, we know that any irreducible component in  $Z_{B_1}$ is pro-modular. In particular, $C$ is pro-modular of dimension $1+2[F:\mathbb{Q}]$.
\end{proof}

\begin{rem}\label{adv}
	1) If $\bar{\chi}$ is quadratic, then we cannot ignore the dihedral condition in the defintion of a nice prime. By Corollary \ref{red-pseudo}, the upper bound of the dihedral locus is $1+[F: \mathbb{Q}]$, so we may not find a nice prime from the pro-modular prime $\mathfrak{s}$ as its dimension may not be large enough. However, we can find one-dimensional irreducible pro-modular prime from $\mathfrak{s} $ and then use Proposition \ref{promodular to nice}, which is the advantage of our method (particularly Corollary \ref{innovation'}) in this paper.
	
	2) We can also prove the same results as Theorem \ref{MINIMAL} using the same proof.
\end{rem}

\newpage
\bibliography{main.bib}
\bibliographystyle{alpha}~
~

Graduate School of Mathematical Sciences, University of Tokyo, Komaba, Meguro, Tokyo 153-8914, Japan\\

\textit{Email address}: zhangxy96@g.ecc.u-tokyo.ac.jp

\end{document}

	%\subsection{Generalized matrix algebras}
	In this subsection, we will briefly introduce generalized matrix algebras, and use GMA as an abbreviation. Most of the details can be found in ,  and \cite{bellaiche2009families}.
	
	Let $R$ be a CNL ring with maximal ideal $\mathfrak{m}_R$ and residue field $\mathbb{F}$, and we view $R$ as a topological ring under the $\mathfrak{m}_R$-adic topology.
	
	\begin{defn}\label{a}
		We say that $A= \begin{pmatrix}
			R & B\\
			C &R
		\end{pmatrix}$ a topological GMA (of type $(1, 1)$) over $R$, if\\
		
		1) $B $ and $C$ are topological $R$-modules,
		
		2) an element of $A$ is of the form
		$\begin{pmatrix}
			a & b\\
			c & d
		\end{pmatrix}$ 
		with $a,d \in R$, $b \in B $ and $c \in C$,
		
		3) a continuous morphism $m': B \otimes_R C \to R$ of $R$-modules is given such that for all $b_1, b_2 \in B$ and $c_1, c_2 \in C$, $m'(b_1 \otimes c_1)b_2 = m'(b_2 \otimes c_1)b_1$ and $m'(b_1 \otimes c_1)c_2 = m'(b_1 \otimes c_2)c_1$.
		
		Further, $A$ is a topological $R$-algebra with the addition given by $$
		\begin{pmatrix}
			a_1 & b_1\\
			c_1 & d_1
		\end{pmatrix}
		+
		\begin{pmatrix}
			a_2 & b_2\\
			c_2 & d_2
		\end{pmatrix}
		=
		\begin{pmatrix}
			a_1+a_2 & b_1+b_2\\
			c_1+ c_2& d_1+d_2
		\end{pmatrix},
		$$
		the multiplication given by$$
		\begin{pmatrix}
			a_1 & b_1\\
			c_1 & d_1
		\end{pmatrix}
		\cdot
		\begin{pmatrix}
			a_2 & b_2\\
			c_2 & d_2
		\end{pmatrix}
		=
		\begin{pmatrix}
			a_1a_2+m'(b_1 \otimes c_2)c_1 & a_1b_2+d_2b_1\\
			d_1c_2+a_2c_1& d_1d_2+m'(b_2 \otimes c_1)
		\end{pmatrix},$$
		and the topology given by those on $R$, $B$ and $C$.
	\end{defn}
	
	From our definition, we can observe that for any topological CNL ring $R$, $A=M_2(R)$ has a unique structure of topological GMA.
	
	\begin{defn}
		Let $A= \begin{pmatrix}
			R & B\\
			C &R
		\end{pmatrix}$ be a topological GMA over a CNL ring $R$. We say that $A'= \begin{pmatrix}
			R & B'\\
			C' &R
		\end{pmatrix}$ is a sub-$R$-GMA of $A$, if $B', C'$ are $R$-submodules of $B,C$ respectively, and the inclusion maps keep the GMA structure.
	\end{defn}
	
	\begin{prop}\label{sub-GMA}
		Let $A= \begin{pmatrix}
			R & B\\
			C &R
		\end{pmatrix}$ be a topological GMA over a CNL ring $R$. Let $\rho: G \to R^\times$ be a representation of a group $G$. Assume that there exists an element $g_0 \in G$ such that $\rho(g_0)$ is diagonal, with diagonal terms distinct modulo $\mathfrak{m}_R$. Then $R[\rho(G)]$ is a sub-$R$-GMA over $A$.
	\end{prop}
	
	\begin{proof}
		See \cite[Lemma 2.4.5]{bellaiche2019image}.
	\end{proof}
	
	Before stating the next proposition, we first give the following setup.
	
	Let $R$ be a CNL ring with residue field $\mathbb{F}$ and $S$ an $R$-algebra (maybe not commutative). Let $A= \begin{pmatrix}
		R & R_{1,2}\\
		R_{2,1} &R
	\end{pmatrix}$ be a topological GMA over $R$ which is also a quotient of $S$. Assume that $m'(R_{1,2} \otimes_{R} R_{2,1}) \subset \mathfrak{m}_R$. Write $e_1=\begin{pmatrix}
		1 & 0\\
		0 & 0
	\end{pmatrix}$ and $e_2=\begin{pmatrix}
		0 & 0\\
		0 & 1
	\end{pmatrix}$ in $A$. Then the map $\rho_i: S/\mathfrak{m}_RS \to A/\mathfrak{m}_RA \to e_i(A/\mathfrak{m}_RA) e_i \cong M_1(\mathbb{F})$ defines a representation of $S/\mathfrak{m}_RS$ for $i=1, 2$. For $i,j \in \{1,2\}$, by an extension of $\rho_j$ by $\rho_i$, we mean a representation $S/\mathfrak{m}_RS \to \operatorname{End}_\mathbb{F} (V)$ together with an exact sequence of $S/\mathfrak{m}_RS$-modules$$
	0 \to \rho_i \to V \to \rho_j \to 0.$$
	Then we know that $V$ is an $\mathbb{F}$-vector space of dimension $2$. Such an extension defines an element in the module $\operatorname{Ext}_{S/\mathfrak{m}_RS}^1 (\rho_j, \rho_i)$.
	
	\begin{prop} \label{GMA} Keep assumptions as above.
		For $i \ne j$, there exists a natural injective $\mathbb{F}$-map: $$
		\iota_{i,j}: \operatorname{Hom}_R (R_{i,j}, \mathbb{F}) \hookrightarrow \operatorname{Ext}_{S/\mathfrak{m}_RS}^1 (\rho_j, \rho_i). $$
	\end{prop}	
	
	\begin{proof}
		It is actually \cite[Theorem 1.5.5]{bellaiche2009families}. We sketch a proof here.
		
		For an element $f \in \operatorname{Hom}_R (R_{i,j}, \mathbb{F})$, we can define an $R$-linear map $\rho_f:  S \to A \to M_2(\mathbb{F})$, 
		$$ x \mapsto \begin{pmatrix}
			a_{i,i} (x)& f(a_{i,j}(x))\\
			0 & a_{j,j}(x)
		\end{pmatrix}, ~x \in S,$$
		where $$a_{i,j}: A/\mathfrak{m}_RA \to e_i(A/\mathfrak{m}_RA) e_j,~ x \mapsto e_ixe_j$$ is the canonical projection. As $m'(R_{1,2} \otimes_{R} R_{2,1}) \subset \mathfrak{m}_R$, we know that $\rho_f$ is an $\mathbb{F}$-algebra homomorphism which defines an extension of $\rho_j$ by $\rho_i$. This defines the map $\iota_{i,j} $, and consequently, the map $\rho_f$ defines an element $\iota_{i,j}(f)$ in $\operatorname{Ext}_{S/\mathfrak{m}_RS}^1 (\rho_j, \rho_i)$.
		
		We only need to show that $\iota_{i,j} $ is actually an injection. Assume that $\iota_{i,j}(f) =0$ for some $f$. Then the extension defined by $f$ is split. From our construction, we know that the map $\rho_f : S \to \operatorname{End}_\mathbb{F} (V)$ factors through $ A/\mathfrak{m}_RA$. Thus, the extension is also split when restricted to any subalgebra of $ A/\mathfrak{m}_RA$. We restrict it to the subalgebra $e_i(A/\mathfrak{m}_RA) e_j$. Then the restricted extension is $$ x \mapsto \begin{pmatrix}
			0& f(a_{i,j}(x))\\
			0 & 0
		\end{pmatrix}, ~x \in S.$$ 
		It is split if and only if $f=0$. This proves our result.
	\end{proof}
	
	\begin{rem}
		If we have a representation $\rho$ and a GMA $A$ satisfying the conditions in Proposition \ref{sub-GMA}, then we can take $S$ as the group algebra $R[G]$ in Proposition \ref{GMA}. As $R[\rho(G)]$ is a GMA over $A$ which is also a quotient of $R[G]$, we can use Proposition \ref{GMA} to study the image $ \rho(G)$. 
	\end{rem}
	
%\section{Applications}

In this section, we use our pro-modularity arguments to study the dimensions of the big Hecke algebra and Galois deformation rings.

%\subsection{The dimension of Hecke algebra} In this subsection, we study the dimension of our big Hecke algebra $(\mathbb{T}_{\xi})_{\mathfrak{m}_{\xi}}$.

In Deo's article \cite{deo2023density}, it is proved that the big Hecke algebra is a reduced local complete intersection ring of Krull dimension $4$. One of the crucial step of his proof is the following \textit{infinite fern} argument.

\begin{thm} \cite[Corollary 2.28]{AST_2011__339__31_0}
	The Krull dimension of each component of the spectrum of the big Hecke algebra for $\mathbb{Q}$ is at least $4$.
\end{thm}

\begin{rem}
	Such a result is first studied by Gouv{\^e}a and Mazur in . The case for totally real fields is studied in 
\end{rem}

Using Pan's patching theorem and our pro-modularity arguments, we can give an analogue result to Deo's work.

\begin{thm}\label{dimT}
	Keep assumptions as in Theorem \ref{pro}. Then $(\mathbb{T}_{\xi})_{\mathfrak{m}_{\xi}}$ is equidimensional of dimension $1+2[F: \mathbb{Q}]$.
\end{thm}

Before giving the proof, we briefly summarize some main steps of the proof of Theorem \ref{R=T}. Here we follow most notations as in \cite[Section 4]{pan2022fontaine}, and one can find more details there.

Let $\mathfrak{q}$ be a nice prime of $(\mathbb{T}_{\xi})_{\mathfrak{m}_{\xi}}$. Let $B$ be the topological closure of the $\mathbb{F}$-algebra generated by all the entries of $\rho(\mathfrak{q})^o(\operatorname{Gal}(F_{\Sigma}/F))$. By Chebotarev’s density Theorem, we may find a finite set of primes $T'$ disjoint with $\Sigma$ such
that the entries of $\rho(\mathfrak{q})^o(\operatorname{Frob}_v), v \in T'$, topologically generate $B$. Let $P= T' \cup \Sigma$.

Let $r=\operatorname{dim}_{k(\mathfrak{q})} H^1(G_{F,P},\operatorname{ad}^0 \rho(\mathfrak{q})(1))$ (the cardinality of Taylor-Wiles primes). Let $R_{\infty}^{\{\xi_v\}}= R_{\textnormal{loc}}^{\{\xi_v\}}[[x_1, ..., x_g]]$, where $ R_{\textnormal{loc}}^{\{\xi_v\}}$ is the completed tensor product of local deformation rings determined by primes in $P$ (see \cite[Definition 4.2.2]{pan2022fontaine}) and $g=r+|P|-[F:\mathbb{Q}]-1$. Let $S_\infty'=\mathcal{O}[[y_1, ... , y_{4|P|-1}, s_1', ... , s_r']]$, and $S_\infty=\mathcal{O}[[y_1, ... , y_{4|P|-1}, s_1, ... , s_r]]$ with an ideal $\mathfrak{a}_1=(y_1, ... , y_{4|P|-1}, s_1, ... , s_r)$ such that $ S_\infty$ is a finite free $S_\infty'$-algebra.

Let $\textnormal{m}_\infty^{\{\xi_v\}}$ and $\textnormal{m}_0^{\{\xi_v\}}$ be the \textit{patched Hecke modules }(see \cite[Section 4.7]{pan2022fontaine}). Then $\textnormal{m}_\infty^{\{\xi_v\}} $ is a flat  $ S_\infty$-module and $\textnormal{m}_\infty^{\{\xi_v\}}/ \mathfrak{a}_1 \textnormal{m}_\infty^{\{\xi_v\}} \cong (\textnormal{m}_0^{\{\xi_v\}})^{\oplus 2^r}$. By Pan's local-global compatibility result \cite[Corollary 3.5.8]{pan2022fontaine}, $\textnormal{m}_0^{\{\xi_v\}}$ is a finitely generated faithful $\widehat{(\mathbb{T}_{\xi})_{\mathfrak{q}}}$-module. By Theorem \ref{dimofT}, we have $\operatorname{dim}_{\widehat{(\mathbb{T}_{\xi})_{\mathfrak{q}}}} (\textnormal{m}_0^{\{\xi_v\}})= \operatorname{dim} \widehat{(\mathbb{T}_{\xi})_{\mathfrak{q}}} \ge 2[F:\mathbb{Q}]$ since $\mathfrak{q}$ is of dimension $1$.

Let $(R^{\{\xi_v\}})'$ be $\widehat{(R_{\infty}^{\{\xi_v\}})_{\mathfrak{q}_{\infty}^{\{\xi_v\}}}} \otimes_{S_\infty'} S_\infty$, where ${\mathfrak{q}}_{\infty}^{\{\xi_v\}} $ is the prime of $ R_{\infty}^{\{\xi_v\}}$ corresponding to the nice prime $\mathfrak{q}$.	We know that $\widehat{(\mathbb{T}_{\xi})_{\mathfrak{q}}} $ is a natural quotient of $(R^{\{\xi_v\}})'$ by mapping $ S_\infty$ to $S_\infty/ \mathfrak{a}_1=\mathcal{O}$. Note that $ y_1, ... , y_{4|P|-1}, s_1, ... , s_r \in \mathfrak{a}_1$ form a regular sequence of $\textnormal{m}_\infty^{\{\xi_v\}}$. Then we can conclude that $$\operatorname{dim}_{(R^{\{\xi_v\}})'}(\textnormal{m}_\infty^{\{\xi_v\}}) \ge 4|P|-1 +r +\operatorname{dim} \widehat{(\mathbb{T}_{\xi})_{\mathfrak{q}}} \ge 4|P|-1 +r+2[F:\mathbb{Q}].$$

Using Taylor's Ihara avoidance trick \cite{PMIHES_2008__108__183_0}, one can show that  
$\widehat{(R_{\infty}^{\{\xi_v\}})_{\mathfrak{q}_{\infty}^{\{\xi_v\}}}}$ is equidimensional of dimension $ 4|P|-1 +r+2[F:\mathbb{Q}]$ (see \cite[Section 4.2, Lemma 4.8.2]{pan2022fontaine}). As $(R^{\{\xi_v\}})'$ is finite free over $\widehat{(R_{\infty}^{\{\xi_v\}})_{\mathfrak{q}_{\infty}^{\{\xi_v\}}}}$, we have $$\operatorname{dim}_{(R^{\{\xi_v\}})'}(\textnormal{m}_\infty^{\{\xi_v\}})=\operatorname{dim}_{\widehat{(R_{\infty}^{\{\xi_v\}})_{\mathfrak{q}_{\infty}^{\{\xi_v\}}}}}(\textnormal{m}_\infty^{\{\xi_v\}})=4|P|-1 +r+2[F:\mathbb{Q}]= \operatorname{dim} \widehat{(R_{\infty}^{\{\xi_v\}})_{\mathfrak{q}_{\infty}^{\{\xi_v\}}}}.$$

If the characters in $\{\xi_v\}$ are all non-trivial, by \cite[Proposition 4.2.3]{pan2022fontaine},  $\widehat{(R_{\infty}^{\{\xi_v\}})_{\mathfrak{q}_{\infty}^{\{\xi_v\}}}}$ is irreducible. From the previous equations, we know that $\textnormal{m}_\infty^{\{\xi_v\}} $ has full support on $\widehat{(R_{\infty}^{\{\xi_v\}})_{\mathfrak{q}_{\infty}^{\{\xi_v\}}}}$. By \cite[Lemma 4.8.4]{pan2022fontaine}, Theorem \ref{R=T} holds in this case. Similar to the proof of \cite[Theorem 4.1]{PMIHES_2008__108__183_0}, one can deduce the general case using Taylor's trick.

From Theorem \ref{dimofT}, \cite[Corollary 12.5]{eisenbud2013commutative}, and discussions above, we can also conclude that for a nice prime $\mathfrak{q}$ of $(\mathbb{T}_{\xi})_{\mathfrak{m}_{\xi}}$, we have $\operatorname{dim} \widehat{(\mathbb{T}_{\xi})_{\mathfrak{q}}} =\operatorname{dim} (\mathbb{T}_{\xi})_{\mathfrak{q}}= 2[F:\mathbb{Q}]$.

\begin{proof}[Proof of Theorem \ref{dimT}]
	Similar to \cite[Corollary 3.3.3]{pan2022fontaine}, we know that $(\mathbb{T}_{\xi})_{\mathfrak{m}_{\xi}}$ is a CNL ring and hence catenary. Then for every irreducible component of $(\mathbb{T}_{\xi})_{\mathfrak{m}_{\xi}}$, if it contains a nice prime, it is of dimension $1+2[F: \mathbb{Q}]$.
	
	From the proof of Theorem \ref{pro}, we know that every irreducible component of $R^{\textnormal{ps}, \{\xi_v\}} $ of dimension at least $1+2[F:\mathbb{Q}]$ contains a nice prime. Combining Theorem \ref{dimofT} and Theorem \ref{R=T}, we know that for every irreducible component of $(\mathbb{T}_{\xi})_{\mathfrak{m}_{\xi}}$ contains a nice prime. This proves our result.
\end{proof}

%\subsection{The dimension of Galois deformation ring} In this subsection, we study the dimension of 
$R_x^{\textnormal{aux}}$.

Keep assumptions as in Theorem \ref{pro}. The following result is an analogue to \cite[Theorem A]{deo2023density} and Proposition \ref{dim aux}.

\begin{thm}\label{dimofR}
	$(1)$ For every irreducible component of $R_\textnormal{aux}^\textnormal{ps} $ of dimension at least $1+2[F:\mathbb{Q}]$, it is of dimension $1+2[F:\mathbb{Q}]$.
	
	$(2)$ As in Theorem \ref{MINIMAL}, we further assume $$[F:\mathbb{Q}] \ge \operatorname{max}\{1+\frac{1}{2}\operatorname{dim}_{\mathbb{F}}(H^1(F_{\Sigma}/F, \mathbb{F}(\bar{\chi}^{-1}))), 7|\Sigma \setminus \Sigma_p|+2\}.$$ 
	Then for any non-zero element $x \in H^1(F_{\Sigma}/F, \mathbb{F}(\bar{\chi}^{-1}))$, $R_x^\textnormal{aux}$ is a local complete intersection ring of Krull dimension $1+2[F:\mathbb{Q}]$.
	
\end{thm}

\begin{proof}
	The first assertion follows from Theorem \ref{R=T} and Theorem \ref{dimT}.
	
	Let $\mathfrak{P}$ be a minimal prime of $ R_x^\textnormal{aux}$. Let $\mathfrak{P}'$ be the inverse image of $\mathfrak{P}$ via the natural map $f: R_\textnormal{aux}^\textnormal{ps} \to R_x^\textnormal{aux}$. From our assumptions, Proposition \ref{raux}, Proposition \ref{red} and Proposition \ref{pan}, we have $$1+2[F:\mathbb{Q}] \le \operatorname{dim} R_x^\textnormal{aux}/\mathfrak{P} \le \operatorname{dim} R_\textnormal{aux}^\textnormal{ps}/\mathfrak{P}'.$$
	Then by the first assertion, we know that $ \mathfrak{P}'$ is of dimension $1+2[F:\mathbb{Q}] $. Consequently, $ R_x^\textnormal{aux}$ is of dimension at most $1+2[F:\mathbb{Q}] $. Using Proposition \ref{dimension} and Proposition \ref{raux}, we obtain that $R_x^\textnormal{aux}$ is a local complete intersection ring of Krull dimension $1+2[F:\mathbb{Q}]$.
	
\end{proof}
%\section*{Appendix: Galois cohomology}

Let $\theta: \operatorname{Gal}(F_{\Sigma}/F) \to \mathbb{F}^{\times}$ be a continuous totally odd character. Write $F_{\theta}$ be the splitting field of the character $\theta \omega_p^{-1}$. Let $\operatorname{Cl}(F_{\theta})$ be the $p$-class group of $F_{\theta}$.

\begin{prop*}
	We have $\operatorname{dim}_{\mathbb{F}}(H^1(F_{\Sigma}/F, \mathbb{F}(\theta^{-1}))) = \operatorname{dim}_{\mathbb{F}} (H^1(F_{\Sigma}/F, \mathbb{F}(\theta)))=[F:\mathbb{Q}]$ if and only if the following condition holds:
	
	1) $\theta \mid_{G_{F_v}} \ne \omega_p^{\pm 1} \mid_{G_{F_v}}$ for all $v \in \Sigma$,
	
	2) $\operatorname{Cl}(F_{\theta}) [\theta \omega_p^{-1}] = \operatorname{Cl}(F_{\theta}) [\theta^{-1} \omega_p] =0$, where $\operatorname{Cl}(F_{\theta}) [\eta] $ is the $\eta$-eigenspace of $\operatorname{Cl}(F_{\theta})$ for any character $\eta: \operatorname{Gal}(F_{\theta}/F) \to \mathbb{F}^{\times}$.
\end{prop*}

\begin{proof}
	It is an analogue to \cite[Lemma 2.21]{deo2023density}. We sketch a proof here.
	
	Let $H_{\Sigma}= \operatorname{ker} (H^1(F_{\Sigma}/F, \mathbb{F}(\theta^{-1} \omega_p)) \to \prod_{v \in \Sigma} H^1(F_v, \mathbb{F}(\theta^{-1} \omega_p)))$ be the "Selmer" group. By the Greenberg-Wiles formula (\cite[Theorem 8.7.9]{neukirch2013cohomology}) and local characteristic formula, we have $$
	\operatorname{dim}_{\mathbb{F}} (H^1(F_{\Sigma}/F, \mathbb{F}(\theta))) = \operatorname{dim}_{\mathbb{F}} (H_\Sigma) + [F:\mathbb{Q}]+ \sum_{v \in \Sigma} \operatorname{dim}_{\mathbb{F}} (H^0(F_v, \mathbb{F}(\theta^{-1} \omega_p))), $$
	where we take the local condition $\mathcal{L}_v$ as $H^1(F_v, \theta\mid_{G_{F_v}}) $. Thus, we have $ \operatorname{dim}_{\mathbb{F}} (H^1(F_{\Sigma}/F, \mathbb{F}(\theta)))=[F:\mathbb{Q}]$ if and only if the following conditions hold:
	
	i) $H^0(F_v, \mathbb{F}(\theta^{-1} \omega_p))=0$ for all $v \in \Sigma$,
	
	ii) the Selmer group $H_{\Sigma}$ is trivial.
	
	Clearly, i) is equivalent to $\theta \mid_{G_{F_v}} \ne \omega_p \mid_{G_{F_v}}$ for all $v \in \Sigma$. 
	
	By inflation-restriction exact sequence, we have $$ H^1 (F_v^{ur}/F_v, (\mathbb{F}(\theta^{-1} \omega_p))^{I_v}) = \operatorname{ker}(H^1(F_v, \mathbb{F}(\theta^{-1} \omega_p)) \to H^1(I_v, \mathbb{F}(\theta^{-1} \omega_p))),~v \in \Sigma,$$
	where $F_v^{ur} $ is the maximal unramified extension of $F_v$. As $ \mathbb{F}$ is a finite field, we have $$ (\mathbb{F}(\theta^{-1} \omega_p))^{I_v}/(\operatorname{Frob}_v-1)(\mathbb{F}(\theta^{-1} \omega_p))^{I_v}=H^1 (F_v^{ur}/F_v, (\mathbb{F}(\theta^{-1} \omega_p))^{I_v}) =0,$$ if $\theta \mid_{G_{F_v}} \ne \omega_p \mid_{G_{F_v}}$. Therefore, $$
	H_{\Sigma}= \operatorname{ker} (H^1(F_{\Sigma}/F, \mathbb{F}(\theta^{-1} \omega_p)) \to \prod_{v \in \Sigma} H^1(I_v, \mathbb{F}(\theta^{-1} \omega_p))).$$
	By class field theory, we know that $ \operatorname{dim}_{\mathbb{F}} (H^1(F_{\Sigma}/F, \mathbb{F}(\theta)))=[F:\mathbb{Q}]$ if and only if the following conditions hold:
	
	i') $\theta \mid_{G_{F_v}} \ne \omega_p \mid_{G_{F_v}}$ for all $v \in \Sigma$,
	
	ii') $\operatorname{Cl}(F_{\theta}) [\theta^{-1} \omega_p] =0 $.
	
	Our result follows from this equivalence.
\end{proof}

\cite{liu2006algebraic}
\cite{Chenevier2011}
